\section{Conventions and exact excitation compression}
\label{supp:conventions}

Let

\begin{equation}\label{eq:supp-labels}
 \mathcal I_n=\{[a,b]:1\le a\le b\le n\},
 \qquad M_n=|\mathcal I_n|=\frac{n(n+1)}2,
\end{equation}

with physical prior \(1/M_n\).  For an overlap \(0\le c_n\le1\), choose
the phase convention

\begin{equation}\label{eq:supp-state-gauge}
 |\psi_n\rangle=c_n|0\rangle+s_n|1\rangle,
 \qquad s_n=\sqrt{1-c_n^2}.
\end{equation}

The interval states have Gram matrix

\begin{equation}\label{eq:supp-physical-gram}
 (G_n)_{I,J}=c_n^{|I\triangle J|}.
\end{equation}

We write

\begin{equation}\label{eq:supp-critical-functions}
 L_n=\log n,
 \qquad
 A(t)=1+\frac{2\sqrt t}{\pi}+\frac{\sqrt2\,t}{\pi^2},
 \qquad F(t)=A(t)^2.
\end{equation}

For a generic overlap \(0\le c<1\), define

\[
 s(c)=\sqrt{1-c^2},\qquad
 J(c)=\frac1{2\pi}\int_0^{2\pi}
 \frac{d\theta}{|1-ce^{i\theta}|},\qquad
 \ell(c)=\pi J(c),
\]

and

\[
 p_1(c)=(1-c^2)J(c)^2
 =\frac{1-c^2}{\pi^2}\ell(c)^2,
 \qquad p_1(1)=0.
\]

For a moving overlap, \(s_n=s(c_n)\) as in
\eqref{eq:supp-state-gauge}, and we put

\[
 \ell_n=\ell(c_n),\qquad
 \lambda_n=n(1-c_n),\qquad
 \tau_n=\lambda_nL_n^2,\qquad
 h_n=np_1(c_n).
\]

For a compact critical parameter \(t\ge0\), set

\begin{equation}\label{eq:supp-critical-family-definition}
 c_n(t)=1-\frac{t}{nL_n^2}.
\end{equation}

Thus \(L_n=\log n\) remains the logarithmic dimension variable, whereas
\(\ell_n\) is the elliptic integral scale.  The functions \(J\) and
\(\ell\) are used only for \(c<1\); the endpoint \(c=1\) is always handled
directly.

For a full \(M_n\)-label Gram matrix, \(P_{\rm tr}\), \(P_{\rm SRM}\), and
\(P_{\rm opt}\) denote the trace benchmark, the square-root measurement
success probability, and the Bayes optimum with external prior \(1/M_n\).
For \(k\in\{0,1,\ge2\}\), \(P_{\rm opt}^{(k)}\) is the optimum for the
subnormalized vectors in excitation sector \(k\), retaining the same external
prior. Singleton interval labels therefore remain present even when their
vectors vanish in the strict-hull sector.

Sections~\ref{supp:trace-laws} and \ref{supp:dual-certificates} close the
compact-window global squeeze through a trace lower bound and sectorwise dual
upper bounds. Section~\ref{supp:matching} extends the squeeze to moving
overlaps and proves the unified hull-dominance law by weighted hull
compression and adaptive endpoint blocks. Section~\ref{supp:direct-srm} is an
independent labelwise SRM refinement. Finally, Section~\ref{supp:numerics}
records the finite-size
numerical diagnostics.

\begin{lemma}[Elliptic endpoint scale]
\label{lem:supp-elliptic-endpoint}
As \(c\uparrow1\),

\begin{equation}\label{eq:supp-elliptic-endpoint}
 \ell(c)=\log\frac{8}{1-c}+o(1).
\end{equation}

Moreover, for every compact \(C\Subset(0,\infty)\),

\begin{equation}\label{eq:supp-elliptic-endpoint-uniform}
 \sup_{\lambda\in C}
 \left|\ell(1-\lambda/n)-\log n-\log(8/\lambda)\right|
 \longrightarrow0.
\end{equation}
\end{lemma}

\begin{proof}
With

\[
 m=\frac{4c}{(1+c)^2},
 \qquad \delta=\sqrt{1-m}=\frac{1-c}{1+c},
\]

the substitution \(\theta=2\phi\) gives

\[
 \ell(c)=\frac{2}{1+c}\,K(m),
\]

where \(K\) uses the parameter convention of the main text. Its endpoint
expansion is

\[
 K(m)=\log\frac4\delta
 +O\!\left(\delta^2\log\frac1\delta\right).
\]

Substitution of \(\delta=(1-c)/(1+c)\) proves
\eqref{eq:supp-elliptic-endpoint}. If \(c=1-\lambda/n\) and \(\lambda\) lies
in a compact subset of \((0,\infty)\), the same remainder is uniform in
\(\lambda\), which proves \eqref{eq:supp-elliptic-endpoint-uniform}.
\end{proof}

When \(c_n>0\), put \(q_n=s_n/c_n\) and
\((D_n)_{I,I}=c_n^{|I|}\).  The state matrix, with rows indexed by interval
labels and columns by excitation subsets, is

\begin{equation}\label{eq:supp-state-matrix}
 (A_n)_{I,S}=c_n^{|I|}q_n^{|S|}\mathbf1_{\{S\subseteq I\}},
 \qquad G_n=A_nA_n^*.
\end{equation}

Thus \(\operatorname{tr}\sqrt{G_n}=\|A_n\|_{S_1}\).  In every bounded
critical window, \(c_n\to1\); hence \(c_n=0\) can occur only in a finite
prefix.  Indices with \(c_n=1\) are handled directly, since
\(G_n=J_{M_n}\) and
\(P_{\rm tr}=P_{\rm SRM}=P_{\rm opt}=M_n^{-1}\).

\begin{proposition}[Exact weighted hull compression]
\label{prop:supp-hull-compression}
Let \(Z_n\) be the proper-interval containment matrix

\begin{equation}\label{eq:supp-zeta}
 (Z_n)_{I,(s,t)}=\mathbf1_{\{s<t,\,[s,t]\subseteq I\}}.
\end{equation}

For \(0<c<1\), set

\[
 |e_c\rangle=\frac{|\psi\rangle-c|0\rangle}{s(c)}
\]

and, for \(1\le u<v\le n\),

\[
 |\eta_{u,v}\rangle
 =|0\rangle^{\otimes(u-1)}|e_c\rangle_u
  |\psi\rangle^{\otimes(v-u-1)}|e_c\rangle_v
  |0\rangle^{\otimes(n-v)}.
\]

These vectors form an orthonormal family, and the exact coordinate state
matrix of the depth-at-least-two sector is

\begin{equation}\label{eq:supp-weighted-hull-identity}
 [A_{\ge2,n}(c)]_{(a,b),(u,v)}
 =(1-c^2)c^{(u-a)+(b-v)}
 \mathbf1_{\{a\le u<v\le b\}}.
\end{equation}

In particular, all \(M_n\) physical rows are retained, while the column set
has cardinality \(n(n-1)/2\). Define

\[
 [V_n(c)]_{a,u}=c^{u-a}\mathbf1_{\{a\le u\}},
\]

and let \(R_{\mathcal I}\) restrict ordered row pairs to \(a\le b\) and
\(C_{<}\) include strict column pairs \(u<v\). Then

\[
 A_{\ge2,n}(c)
 =(1-c^2)R_{\mathcal I}
 [V_n(c)\otimes V_n(c)^{\mathsf T}]C_{<}.
\]

Equivalently, with

\[
 u^+_{a,c}(x)=s(c)c^{x-a}\mathbf1_{\{x\ge a\}},\qquad
 u^-_{b,c}(y)=s(c)c^{b-y}\mathbf1_{\{y\le b\}},
\]

the row belonging to \(I=[a,b]\) is
\(P_{\{x<y\}}(u^+_{a,c}\otimes u^-_{b,c})\).

For the moving critical notation \(0<c_n<1\), the same matrix is, exactly,

\begin{equation}\label{eq:supp-exact-hull}
 A_{\ge2,n}=q_n^2D_nZ_nW_n,
 \qquad (W_n)_{(s,t),(s,t)}=c_n^{-(t-s-1)}.
\end{equation}

Consequently,

\begin{equation}\label{eq:supp-hull-order}
 I\preceq W_n\preceq\rho_n I,
 \qquad \rho_n=c_n^{-(n-2)}.
\end{equation}

Define the comparison Gram matrix by

\begin{equation}\label{eq:supp-unweighted-comparison-gram}
 \widetilde G_n
 =D_n\bigl(J_{M_n}+q_n^2B_{1,n}B_{1,n}^*
 +q_n^4Z_nZ_n^*\bigr)D_n,
\end{equation}

where \([B_{1,n}]_{I,x}=\mathbf1_{\{x\in I\}}\) is the interval-site
incidence matrix. This retains the physical length factor \(D_n\) while
replacing the strict-hull depth weights by one. Then

\begin{equation}\label{eq:supp-full-hull-sandwich}
 \widetilde G_n\preceq G_n\preceq\rho_n^2\widetilde G_n.
\end{equation}
\end{proposition}

\begin{proof}
First suppose \(0<c<1\).  Expanding
\(|\psi\rangle=c|0\rangle+s(c)|e_c\rangle\), every excitation subset
\(S\subseteq[a,b]\) with \(|S|\ge2\) has unique first and last sites
\(u=\min S<v=\max S\).  Write
\(S=\{u,v\}\cup T\), where
\(T\subseteq\{u+1,\ldots,v-1\}\).  For a fixed hull \(u<v\), summing the
unrestricted interior sites back to
\(|\psi\rangle^{\otimes(v-u-1)}\) leaves the scalar

\[
 s(c)^2c^{|[a,b]|-(v-u-1)-2}
 =s(c)^2c^{(u-a)+(b-v)}.
\]

This proves \eqref{eq:supp-weighted-hull-identity} and also proves that the
whole projected state lies in the displayed hull span.  If two strict pairs
have different first sites, the smaller first site carries \(|e_c\rangle\)
in one vector and \(|0\rangle\) in the other.  If their first sites agree
but their last sites differ, the larger last site gives the same
\(|e_c\rangle\)-versus-\(|0\rangle\) distinction.  Hence the hull vectors
are orthonormal.  The projected tensor formula now follows entrywise from

\[
 [V_n(c)]_{a,u}[V_n(c)]_{v,b}
 =c^{(u-a)+(b-v)}
  \mathbf1_{\{a\le u\}}\mathbf1_{\{v\le b\}}.
\]

For \(c=c_n\in(0,1)\), direct simplification gives

\[
 q_n^2c_n^{|I|}c_n^{-(v-u-1)}
 =s_n^2c_n^{(u-a)+(b-v)},
\]

Hence \eqref{eq:supp-exact-hull} is the exact entrywise representation of the
physical hull sector in \eqref{eq:supp-weighted-hull-identity}.
The bounds \eqref{eq:supp-hull-order} and
\eqref{eq:supp-full-hull-sandwich} then follow as before, because the
vacuum and one-excitation contributions are unchanged.

For a complex original overlap
\(\langle0|\psi\rangle=ce^{i\varphi}\), replacing
\(|\psi\rangle\) by \(e^{-i\varphi}|\psi\rangle\) multiplies the interval
vector \(I\) by \(e^{-i|I|\varphi}\).  The Gram matrix changes only by
diagonal-unitary congruence, and the physical rank-one hypotheses are
unchanged.  This justifies the nonnegative phase gauge used above.

At \(c=1\), every hypothesis is the vacuum and the hull sector vanishes, so
neither \(|e_c\rangle\) nor \(q_n\) is needed. At \(c=0\), set
\(|e_0\rangle=|\psi\rangle\). A proper interval then has the single hull column
\((u,v)=(a,b)\), while singleton rows vanish; equation
\eqref{eq:supp-weighted-hull-identity} remains valid with \(0^0=1\) at zero
boundary depth. In both endpoint cases, all \(M_n\) physical labels, including
singleton intervals, retain the external prior \(1/M_n\).
\end{proof}

\begin{lemma}[Subnormalized PGM and dual bounds]
\label{lem:supp-subnormalized-pgm}
For possibly subnormalized equal-prior vectors \(v_1,\ldots,v_M\) with Gram
matrix \(X\), define

\[
 P_{\rm PGM}(v):=\frac1M\sum_{j=1}^M(\sqrt X)_{jj}^2.
\]

Then

\begin{equation}\label{eq:supp-subnormalized-pgm}
 \left(\frac{\operatorname{tr}\sqrt X}{M}\right)^2
 \le P_{\rm PGM}(v)\le P_{\rm opt}(v).
\end{equation}

Moreover, if \(K\succeq|v_j\rangle\langle v_j|\) for all \(j\), then
\(MP_{\rm opt}(v)\le\operatorname{tr}K\).
\end{lemma}

\begin{proof}
The polar measurement gives
\(P_{\rm PGM}=M^{-1}\sum_j(\sqrt X)_{jj}^2\), including singular \(X\) by
support restriction.  Cauchy--Schwarz proves the first inequality.  The
operator \(K/M\) satisfies every equal-prior Bayes-dual constraint, proving
the last assertion.
\end{proof}

\section{Følner and transfer proofs}
\label{supp:folner-transfer-proofs}

This section proves the local diagonal Følner lemma and the exceptional-sector
transfer theorem stated in the main text.

\subsection{Proof of the local diagonal Følner lemma}

We retain the notation and assumptions of the local diagonal Følner lemma in
the main text.

\begin{proof}[Proof of the local diagonal Følner lemma]
Since \(A\) has spectrum in \([\mu,L]\), the same is true of every compression
\(A_n\). Fix \(\varepsilon>0\) and put \(\eta=\varepsilon/8\). By the
Weierstrass theorem, applied to the real and imaginary parts if necessary,
choose a polynomial

\[
p(t)=\sum_{m=0}^{s}a_mt^m
\quad\text{such that}\quad
\sup_{t\in[\mu,L]}|\phi(t)-p(t)|<\eta.
\]

Let \(k_R=k\,\mathbf 1_{\{\|u\|_\infty\le R\}}\), let \(A_R\) be convolution
by \(k_R\), and set \(A_{R,n}=P_nA_RP_n|_{\ell^2(\Lambda_n)}\). Young's
inequality gives

\begin{equation}\label{eq:folner-young}
\|A-A_R\|_{\rm op},\ \|A_n-A_{R,n}\|_{\rm op}
\le\|k-k_R\|_{\ell^1}.
\end{equation}

Choose \(R\) so large that the right hand side of
\eqref{eq:folner-young} is at most \(1\) and

\[
C_p\|k-k_R\|_{\ell^1}<\eta,\qquad
C_p:=\sum_{m=1}^{s}m|a_m|(L+1)^{m-1},
\]

where the second condition is void if \(C_p=0\). The telescoping identity for
\(X^m-Y^m\), together with
\(\|A\|,\|A_n\|\le L\) and
\(\|A_R\|,\|A_{R,n}\|\le L+1\), then yields

\begin{equation}\label{eq:folner-polynomial-truncation}
\|p(A)-p(A_R)\|_{\rm op},
\ \|p(A_n)-p(A_{R,n})\|_{\rm op}<\eta.
\end{equation}

Define

\[
Q_t=\{h\in\mathbb Z^d:\|h\|_\infty\le t\}
\]

and the \(sR\)-interior and boundary by

\begin{align*}
\Lambda_n^{\circ,sR}
&=\{x\in\Lambda_n:x+Q_{sR}\subseteq\Lambda_n\},\\
\partial_{sR}\Lambda_n
&=\Lambda_n\setminus\Lambda_n^{\circ,sR}.
\end{align*}

Because a polynomial of degree \(s\) in an operator of propagation range
\(R\) has propagation radius \(sR\), every path of length \(m\le s\) generated
by the kernel stays in \(\Lambda_n\) when it starts at
\(x\in\Lambda_n^{\circ,sR}\). Therefore

\begin{equation}\label{eq:folner-interior-equality}
\langle\delta_x,p(A_{R,n})\delta_x\rangle
=\langle\delta_0,p(A_R)\delta_0\rangle,
\qquad x\in\Lambda_n^{\circ,sR}.
\end{equation}

Moreover,

\begin{equation}\label{eq:folner-boundary-inclusion}
\partial_{sR}\Lambda_n
\subseteq
\bigcup_{\|h\|_\infty\le sR}
\bigl(\Lambda_n\setminus(\Lambda_n-h)\bigr).
\end{equation}

The index set in this union is finite. Each set has
\(o(|\Lambda_n|)\) elements by the Følner property, and hence
\(|\partial_{sR}\Lambda_n|/|\Lambda_n|\to0\).
With
\(K_p=\sup_{|t|\le L+1}|p(t)|\), define

\begin{align*}
d_{n,R,p}(x)
&:=
\left|
\langle\delta_x,p(A_{R,n})\delta_x\rangle
-\langle\delta_0,p(A_R)\delta_0\rangle
\right|,\\
D_{n,R,p}
&:=\frac1{|\Lambda_n|}\sum_{x\in\Lambda_n}d_{n,R,p}(x).
\end{align*}

Equations
\eqref{eq:folner-interior-equality}--\eqref{eq:folner-boundary-inclusion} give

\begin{equation}\label{eq:folner-boundary-average}
D_{n,R,p}
\le2K_p\frac{|\partial_{sR}\Lambda_n|}{|\Lambda_n|}
\longrightarrow0.
\end{equation}

Finally, the spectral theorem and Fourier transformation identify

\[
\langle\delta_0,\phi(A)\delta_0\rangle
=\frac1{(2\pi)^d}\int_{[0,2\pi]^d}\phi(F(\theta))\,d\theta
=a_\phi.
\]

Put

\[
D_n^\phi:=\frac1{|\Lambda_n|}\sum_{x\in\Lambda_n}
\left|\langle\delta_x,\phi(A_n)\delta_x\rangle-a_\phi\right|.
\]

For all sufficiently large \(n\), functional calculus,
\eqref{eq:folner-polynomial-truncation}, and
\eqref{eq:folner-boundary-average} yield

\begin{align*}
D_n^\phi
&\le
\|\phi(A_n)-p(A_n)\|_{\rm op}\\
&\quad+\|p(A_n)-p(A_{R,n})\|_{\rm op}\\
&\quad+D_{n,R,p}\\
&\quad+\|p(A_R)-p(A)\|_{\rm op}\\
&\quad+\|p(A)-\phi(A)\|_{\rm op}\\
&<4\eta+\frac{\varepsilon}{2}=\varepsilon.
\end{align*}

This proves the claimed local diagonal limit.
\end{proof}

\subsection{Proof of the exceptional-sector transfer theorem}

We retain the notation and assumptions of the exceptional-sector transfer
theorem in the main text.

\begin{proof}[Proof of the transfer theorem]
Put \(d_{n,j}=(\sqrt{H_n})_{jj}\) and
\(\bar d_n=m_n^{-1}\sum_jd_{n,j}\). Positivity and the unit diagonal give
\(0\le d_{n,j}\le1\), because

\[
d_{n,j}^2
\le\langle e_j,H_ne_j\rangle=1.
\]

The local-diagonal hypothesis implies
\(\bar d_n\to\alpha\) and

\[
\frac1{m_n}\sum_jd_{n,j}^2\longrightarrow\alpha^2.
\]

The latter is \(P_{\rm SRM}(H_n)\). Put

\[
a_n=\frac1{m_n}\sum_j|d_{n,j}-\alpha|.
\]

Then \(|\bar d_n-\alpha|\le a_n\to0\), so \(\bar d_n\ge\alpha/2\) for all
sufficiently large \(n\), and

\[
\frac1{2m_n\bar d_n}\sum_j|d_{n,j}-\bar d_n|
\le\frac{a_n}{\bar d_n}\longrightarrow0.
\]

The singular-safe dual certificate proved in the main text now gives

\[
P_{\rm opt}(H_n)\longrightarrow\alpha^2.
\]

By the operator-norm Gram stability lemma and the retained-family
comparison hypothesis, the retained physical family has the same limits:

\begin{equation}
\label{eq:transfer-retained-limits}
P_{\rm opt}(G_{n,\mathcal B}),
P_{\rm SRM}(G_{n,\mathcal B})\longrightarrow\alpha^2.
\end{equation}

The same square root perturbation bound also yields

\begin{equation}
\label{eq:transfer-retained-trace}
\frac{\operatorname{tr}\sqrt{G_{n,\mathcal B}}}{m_n}
\longrightarrow\alpha.
\end{equation}

Write \(\beta_n=m_n/M_n\). Restricting a full ensemble POVM to the retained
labels and completing the resulting sub-POVM gives the upper bound below;
extending an optimal retained family POVM gives the lower bound:

\begin{equation}
\label{eq:transfer-optimum-sandwich}
\beta_nP_{\rm opt}(G_{n,\mathcal B})
\le P_{\rm opt}(G_n)
\le\beta_nP_{\rm opt}(G_{n,\mathcal B})+(1-\beta_n).
\end{equation}

The retained-mass hypothesis, \eqref{eq:transfer-retained-limits}, and
\eqref{eq:transfer-optimum-sandwich} prove
\(P_{\rm opt}(G_n)\to\alpha^2\).

It remains to transfer the SRM. Put

\[
t_n=\frac{\operatorname{tr}\sqrt{G_n}}{M_n}.
\]

Eigenvalue interlacing for the principal submatrix gives

\[
t_n\ge
\beta_n\frac{\operatorname{tr}\sqrt{G_{n,\mathcal B}}}{m_n}
\longrightarrow\alpha.
\]

On the other hand, \(t_n^2\le P_{\rm opt}(G_n)\to\alpha^2\), so
\(t_n\to\alpha\). The universal squeeze

\[
t_n^2\le P_{\rm SRM}(G_n)\le P_{\rm opt}(G_n)
\]

proves the stated SRM conclusion.
\end{proof}

\section{Retained orthogonality}
\label{supp:retained-orthogonality}

The global trace argument and the local refinement use the same retained
subspaces, which we now construct.

Define the unweighted three-stratum comparison

\begin{equation}\label{eq:supp-unweighted-three-stratum}
 \widehat G_n=J_{M_n}+q_n^2B_{1,n}B_{1,n}^*
 +q_n^4Z_nZ_n^*,
\end{equation}

where \(Z_n\) retains its zero singleton rows.  Let
\(z_n=B_{1,n}^*\mathbf1\), and let \(Q_{1,n}\) project onto \(z_n^\perp\).
Then \(\operatorname{ran}(B_{1,n}Q_{1,n})\perp\mathbf1\).

Partition the sites into
\(r_n=\lfloor\sqrt{\log(n+1)}\rfloor\) balanced consecutive blocks
\(E_1,\ldots,E_{r_n}\).  For a left block \(E_p\), set

\[
 F_p(a,s)=\mathbf1_{\{a\le s\}},\qquad s\in E_p,
\]

and let \(Q_p^L\) project onto the orthogonal complement of
\(\operatorname{span}\{\mathbf1_{E_p},F_p^*\mathbf1_n\}\).  For a right
block \(E_r\), define
\(R_r(b,t)=\mathbf1_{\{t\le b\}}\) and let \(Q_r^R\) project onto the
orthogonal complement of
\(\operatorname{span}\{\mathbf1_{E_r},R_r^*\mathbf1_n\}\).  On pair
columns \(E_p\times E_r\) with \(p<r\), retain
\(Q_p^L\otimes Q_r^R\); retain zero when both sites lie in one block.  Their
direct sum is denoted by \(Q_{2,n}\).

\begin{lemma}[Exact retained support and row orthogonality]
\label{lem:supp-retained-orthogonality}
For every \(p<r\), the \((p,r)\) column block of \(Z_nQ_{2,n}\) is
supported on interval rows whose endpoint pair belongs to
\(E_p\times E_r\).  On that cell it factorizes, up to reversal of the right
coordinate, as

\begin{equation}\label{eq:supp-retained-cell-factorization}
 (U_{|E_p|}Q_p^L)\otimes(U_{|E_r|}Q_r^R).
\end{equation}

Different cells have orthogonal row and column supports, and the full
retained row range is orthogonal to every endpoint-additive function
\(f(a)+g(b)\).
\end{lemma}

\begin{proof}
If \(a\) lies before \(E_p\), its row in \(F_p\) is constant and is killed
by \(Q_p^L\); if it lies after \(E_p\), that row is zero.  Thus
\(F_pQ_p^L\) is supported on \(a\in E_p\).  Moreover,
\(\mathbf1_n^*F_pQ_p^L=0\) by the definition of \(Q_p^L\).  The analogous
statements hold for \(R_rQ_r^R\).  Since \(p<r\), the Cartesian support cell
lies entirely in \(a<b\), where containment factorizes, proving
\eqref{eq:supp-retained-cell-factorization}.  Every tensor output has zero
sum in each endpoint coordinate, so it is orthogonal to \(f(a)+g(b)\).
Distinct block pairs have disjoint row and column supports.
\end{proof}

\section{Incidence trace laws and the global trace theorem}
\label{supp:trace-laws}

For the interval-site incidence matrix \(B_{1,n}\) defined above, a direct
count gives

\begin{equation}\label{eq:supp-b1-normal}
 B_{1,n}^*B_{1,n}=(n+1)L_D^{-1},
\end{equation}

where \(L_D\) is the Dirichlet path Laplacian.  Therefore

\begin{equation}\label{eq:supp-b1-singular-values}
 \sigma_k(B_{1,n})=
 \frac{\sqrt{n+1}}{2\sin[\pi k/(2(n+1))]},
 \qquad1\le k\le n,
\end{equation}

and endpoint harmonic summation yields

\begin{equation}\label{eq:supp-b1-trace-law}
 \|B_{1,n}\|_{S_1}
 =\frac{n^{3/2}}\pi\log n+O(n^{3/2}),
 \qquad \|B_{1,n}\|_{\rm op}=O(n^{3/2}).
\end{equation}

The corresponding two-endpoint incidence operator is compressed to a
triangular domain rather than a square one.

\begin{proposition}[Triangular Volterra trace law]
\label{prop:supp-b2-trace}
For the matrix \(Z_n\) in \eqref{eq:supp-zeta},

\begin{equation}\label{eq:supp-b2-trace-law}
 \|Z_n\|_{S_1}
 \sim\frac{n^2(\log n)^2}{2\pi^2}.
\end{equation}

More quantitatively, for an integer \(2\le R\le n\), partition the sites into
\(R\) balanced consecutive blocks of sizes \(\ell_p\). Retain only matched row
and column cells with
\(p<r\), and project each one-dimensional column factor off the constant and
incidence-sum directions used in
Lemma~\ref{lem:supp-retained-orthogonality}. The resulting rectangular
pinching \(Z_{n,R}^{\rm ret}\) satisfies

\begin{align}
 \|Z_{n,R}^{\rm ret}\|_{S_1}
 &\ge\frac{n^2(\log n)^2}{2\pi^2}
 \left[1-C\left(
 \frac1R+\frac{\log R}{\log n}
 +\frac1{\log n}+\frac Rn\right)\right].
 \label{eq:supp-b2-retained-rate}
\end{align}
\end{proposition}

\begin{proof}
For \(U_m(i,j)=\mathbf1_{\{i\le j\}}\),

\begin{equation}\label{eq:supp-volterra-trace}
 \|U_m\|_{S_1}=\frac m\pi\log m+O(m).
\end{equation}

On a matched off-diagonal cell \((p,r)\), the unprojected block is
\(U_{\ell_p}\otimes U_{\ell_r}^{\mathsf T}\).  Each projection removes at
most two column directions, and therefore

\[
 \|U_\ell(I-Q)\|_{S_1}
 \le\sqrt2\,\|U_\ell\|_{S_2}=O(\ell).
\]

Thus each one-dimensional trace norm changes by at most \(O(\ell_p)\) or
\(O(\ell_r)\).  The retained row and column cells are mutually orthogonal,
so rectangular pinching and \eqref{eq:supp-volterra-trace} give

\[
 \|Z_{n,R}^{\rm ret}\|_{S_1}
 \ge\frac1{\pi^2}\sum_{p<r}\ell_p\ell_r
 (\log\ell_p-C)(\log\ell_r-C).
\]

Using
\(\sum_{p<r}\ell_p\ell_r=(n^2-\sum_p\ell_p^2)/2\) and balancedness proves
\eqref{eq:supp-b2-retained-rate}.

For the full upper bound, fix an integer \(R\ge2\) and decompose both endpoint triangles
into all macroblocks.  Blocks with two matching endpoint indices are
compressions of \(U\otimes U\); blocks with one matching index are
compressions of \(U\otimes J\); and blocks with no matching index have rank
one.  For a balanced side length \(\ell\),
\[
 \|U_\ell\otimes J_\ell\|_{S_1}
 =O(\ell^2\log\ell),
 \qquad
 \|J_\ell\otimes J_\ell\|_{S_1}=O(\ell^2).
\]
Thus, after summing the finitely many macroblocks for fixed \(R\), the blocks
with one or zero matching endpoint indices contribute
\(o(n^2\log^2n)\).  The off-diagonal cells with two matching endpoint indices
have total normalized contribution

\[
 \frac1{\pi^2n^2\log^2n}
 \sum_{p<r}\ell_p\ell_r\log\ell_p\log\ell_r
 \longrightarrow\frac{1-R^{-1}}{2\pi^2}.
\]

The \(R\) diagonal cells are rescaled triangular problems.  Their total
limsup contribution is at most \(C^*/R\), because
\(R(n/R)^2\log^2(n/R)/[n^2\log^2n]\to R^{-1}\).  To justify the use of the
finite constant \(C^*\), note first that, up to the fixed reversal of the
right-endpoint coordinate,
\[
 \Lambda_m=\{(x,y)\in\{1,\ldots,m\}^2:x+y\le m+1\},
 \qquad
 Z_n=P_{\Lambda_{n-1}}
 (U_{n-1}\otimes U_{n-1})P_{\Lambda_{n-1}}.
\]
The ideal property and \eqref{eq:supp-volterra-trace} therefore give
\[
 \|Z_n\|_{S_1}
 \le\|U_{n-1}\|_{S_1}^2
 =O(n^2\log^2n).
\]
Hence, if
\(C^*=\limsup_n\|Z_n\|_{S_1}/[n^2\log^2n]\), the block triangle inequality
gives

\[
 C^*\le\frac{1-R^{-1}}{2\pi^2}+\frac{C^*}{R}.
\]

Thus \(C^*\le1/(2\pi^2)\).  The retained lower bound followed by
\(R\to\infty\) gives the reverse inequality and proves
\eqref{eq:supp-b2-trace-law}.
\end{proof}

These sector estimates yield a global trace law without assuming local
square-root additivity.

\begin{theorem}[Critical global square-root trace]
\label{thm:supp-global-trace}
Let \(0\le\tau_n=n(1-c_n)(\log n)^2\le T\).  Then

\begin{equation}\label{eq:supp-global-trace-rate}
 \frac{\operatorname{tr}\sqrt{G_n}}{\sqrt{M_n}}
 \ge A(\tau_n)-C_T\frac{\log\log n}{\log n},
\end{equation}

and, uniformly for \(\tau_n\in[0,T]\),

\begin{equation}\label{eq:supp-global-trace-limit}
 \frac{\operatorname{tr}\sqrt{G_n}}{\sqrt{M_n}}
 -A(\tau_n)\longrightarrow0.
\end{equation}
\end{theorem}

\begin{proof}
Choose
\(R_n=\lfloor\log n/\log\log n\rfloor\).  Work first with the unweighted
state matrix

\[
 \widehat A_n=[\,\mathbf1,\ q_nB_{1,n},\ q_n^2Z_n\,].
\]

Project the one-excitation columns by \(Q_{1,n}\), and use the block
projections of Lemma~\ref{lem:supp-retained-orthogonality} in the
two-endpoint sector.  The three retained row ranges are mutually orthogonal,
so the trace norm of their concatenation is the sum of their trace norms.
The one-excitation loss is explicitly bounded by

\begin{equation}\label{eq:supp-b1-retained-loss}
 \frac{q_n}{\sqrt{M_n}}
 \bigl(\|B_{1,n}\|_{S_1}-\|B_{1,n}Q_{1,n}\|_{S_1}\bigr)
 \le\frac{q_n\|B_{1,n}\|_{\rm op}}{\sqrt{M_n}}
 =O_T((\log n)^{-1}).
\end{equation}

For the two-endpoint blocks, the finite-codimension bound in the proof of
Proposition~\ref{prop:supp-b2-trace} is \(O(\ell)\) per Volterra factor,
hence a relative \(O(1/\log\ell)=O(1/\log n)\) loss.  Consequently,
\eqref{eq:supp-b2-retained-rate} applies to these same projections.

Equations \eqref{eq:supp-b1-trace-law} and
\eqref{eq:supp-b2-retained-rate}, together with

\begin{align*}
 q_n\sqrt n\log n
 &=\sqrt{2\tau_n}\left[
 1+O_T\!\left((n(\log n)^2)^{-1}\right)\right],\\
 q_n^2n(\log n)^2
 &=2\tau_n\left[
 1+O_T\!\left((n(\log n)^2)^{-1}\right)\right]
\end{align*}

give the unweighted version of \eqref{eq:supp-global-trace-rate}.  Indeed,
the error in \eqref{eq:supp-b2-retained-rate} is

\[
 O\!\left(R_n^{-1}+\frac{\log R_n}{\log n}
 +\frac1{\log n}+\frac{R_n}{n}\right)
 =O\!\left(\frac{\log\log n}{\log n}\right).
\]

In particular, for the unweighted comparison matrix in
\eqref{eq:supp-unweighted-three-stratum},

\begin{equation}\label{eq:supp-unweighted-global-trace-limit}
 \sup_{0\le t\le T}\left|
 \frac{\operatorname{tr}\sqrt{\widehat G_n}}{\sqrt{M_n}}-A(t)
 \right|\longrightarrow0.
\end{equation}

To restore the physical weights, observe that
\(\|D_n-I\|_{\rm op}=O_T((\log n)^{-2})\) and
\(\|W_n-I\|_{\rm op}=O_T((\log n)^{-2})\). The ideal property gives

\begin{align*}
 \|(D_n-I)\widehat A_n\|_{S_1}
 &\le\|D_n-I\|_{\rm op}\|\widehat A_n\|_{S_1},\\
 q_n^2\|D_nZ_n(W_n-I)\|_{S_1}
 &\le q_n^2\|W_n-I\|_{\rm op}\|Z_n\|_{S_1}.
\end{align*}

After division by \(\sqrt{M_n}\), both errors are
\(O_T((\log n)^{-2})\). The trace norm is Lipschitz under these perturbations
of the state matrix, so the physical matrix obeys
\eqref{eq:supp-global-trace-rate}.

For the matching upper limit, apply the triangle inequality for the Schatten
\(1\)-norm to
the vacuum, one-excitation, and exact hull sectors, then use
\eqref{eq:supp-b1-trace-law}, \eqref{eq:supp-b2-trace-law}, and
Proposition~\ref{prop:supp-hull-compression}.  This gives
uniformly on compact parameter windows,
\[
 \limsup_{n\to\infty}
 \left[
  \frac{\operatorname{tr}\sqrt{G_n}}{\sqrt{M_n}}-A(\tau_n)
 \right]
 \le 0.
\]
Combined with the retained lower bound, this proves
\eqref{eq:supp-global-trace-limit}.
\end{proof}

The sector upper bound is proved in
Section~\ref{supp:dual-certificates}, Theorem~\ref{thm:supp-sector-upper}.
Together with Theorem~\ref{thm:supp-global-trace} and
\(P_{\rm tr}\le P_{\rm SRM}\le P_{\rm opt}\), it gives the common rate on
compact parameter windows. Retained rectangular pinching supplies the lower
estimate, and dual feasibility supplies the upper estimate.

\section{Cyclic and Volterra dual certificates}
\label{supp:dual-certificates}

For vectors \(v_1,\ldots,v_R\), define the geometric dual value

\begin{equation}\label{eq:supp-geometric-dual}
 \kappa(v_1,\ldots,v_R)=
 \inf\{\operatorname{tr}K:K\succeq0,\ K\succeq|v_j\rangle\langle v_j|
 \text{ for every }j\}.
\end{equation}

The singular-safe rank-one criterion used below is

\begin{equation}\label{eq:supp-rank-one-criterion}
 |v\rangle\langle v|\preceq K
 \quad\Longleftrightarrow\quad
 v\in\operatorname{ran}K\ \text{ and }\
 \langle v,K^+v\rangle\le1.
\end{equation}

\subsection{The cyclic one-excitation certificate}

Let \(u_{a,L}\) be the indicator of the oriented length-\(L\) arc beginning
at \(a\) on the cycle \(\mathbb Z_n\), and set

\begin{equation}\label{eq:supp-cyclic-frame}
 \Sigma_n=\sum_{a\in\mathbb Z_n}\sum_{L=1}^n
 |u_{a,L}\rangle\langle u_{a,L}|.
\end{equation}

\begin{lemma}[Cyclic spectrum and leverage]
\label{lem:supp-cyclic-spectrum}
In the normalized Fourier basis, the eigenvalues of \(\Sigma_n\) are

\begin{equation}\label{eq:supp-cyclic-eigenvalues}
 \lambda_0=\frac{n(n+1)(2n+1)}6,
 \qquad
 \lambda_k=\frac{n}{2\sin^2(\pi k/n)},\quad1\le k<n.
\end{equation}

Consequently,

\begin{align}
 \operatorname{tr}\sqrt{\Sigma_n}
 &=\frac{\sqrt2}{\pi}n^{3/2}\log n+O(n^{3/2}),
 \label{eq:supp-cyclic-root-trace}\\
 \max_{1\le L\le n}
 \langle u_{0,L},\Sigma_n^{-1/2}u_{0,L}\rangle
 &\le\frac{\sqrt2}{\pi}\frac{\log n}{\sqrt n}
 +O(n^{-1/2}).
 \label{eq:supp-cyclic-leverage}
\end{align}
\end{lemma}

\begin{proof}
Translation invariance diagonalizes \(\Sigma_n\).  For \(k\ne0\),

\[
 \sum_{L=1}^n\left|\sum_{x=0}^{L-1}e^{2\pi ikx/n}\right|^2
 =\frac1{\sin^2(\pi k/n)}
 \sum_{L=1}^n\sin^2(\pi kL/n)
 =\frac{n}{2\sin^2(\pi k/n)},
\]

while the zero mode gives \(\sum_LL^2\).  Endpoint harmonic summation proves
\eqref{eq:supp-cyclic-root-trace}.  The leverage is

\[
 \frac{L^2}{n\sqrt{\lambda_0}}
 +\frac{\sqrt2}{n^{3/2}}
 \sum_{k=1}^{n-1}
 \frac{\sin^2(\pi kL/n)}{\sin(\pi k/n)}.
\]

Use \(\sin^2x=(1-\cos2x)/2\), the endpoint expansion of the cosecant, and
the uniform lower bound for partial cosine-harmonic sums
\(\inf_{N,x}\sum_{k=1}^N\cos(kx)/k> -\infty\).  The sum is at most
\((n/\pi)\log n+O(n)\), uniformly in \(L\), proving
\eqref{eq:supp-cyclic-leverage}.
\end{proof}

Let \(t_n\) denote the maximum leverage in
\eqref{eq:supp-cyclic-leverage}.  By
\eqref{eq:supp-rank-one-criterion},

\begin{equation}\label{eq:supp-cyclic-certificate}
 K_n^{\rm cyc}=t_n\sqrt{\Sigma_n}
 \succeq|u_{a,L}\rangle\langle u_{a,L}|,
 \qquad
 \operatorname{tr}K_n^{\rm cyc}
 \le\frac{2}{\pi^2}n(\log n)^2+Cn\log n.
\end{equation}

Every linear interval is a cyclic arc.  Its physical one-excitation vector is
\(s_nc_n^{L-1}u_{a,L}\), whose scalar factor is at most \(s_n\).  Thus the
physical one-excitation dual is \(s_n^2K_n^{\rm cyc}\).

\subsection{Endpoint odd-cosine estimates and the boundary repair}

Let \(r_{m,j}=\mathbf1_{\{j,\ldots,m\}}\), let \(U_m\) have these row
vectors, and set

\begin{equation}\label{eq:supp-volterra-definitions}
 S_m=U_m^{\mathsf T}U_m,
 \qquad C_m=S_m^{1/2},
 \qquad
 \gamma_{m,j}=\langle r_{m,j},S_m^{-1/2}r_{m,j}\rangle.
\end{equation}

\begin{lemma}[Odd-cosine endpoint bounds]
\label{lem:supp-odd-cosine}
There is an absolute \(C\) such that

\begin{equation}\label{eq:supp-leverage-profile}
 \gamma_{m,j}\le\frac1\pi
 \left(\log m+\log_+\frac mj\right)+C,
 \qquad1\le j\le m,
\end{equation}

and

\begin{align}
 \operatorname{tr}C_m&=\frac m\pi\log m+O(m),
 \label{eq:supp-cm-trace}\\
 \sum_{x=1}^m\left|(C_m)_{xx}
 -\frac{\operatorname{tr}C_m}{m}\right|&=O(m).
 \label{eq:supp-cm-flatness}
\end{align}
\end{lemma}

\begin{proof}
Reversing coordinates turns \(U_mU_m^{\mathsf T}\) into the Brownian
covariance matrix.  With

\[
 \theta_k=\frac{(2k-1)\pi}{2m+1},
 \qquad
 \lambda_k=\frac1{4\sin^2(\theta_k/2)},
 \qquad
 v_k(p)=\frac{2}{\sqrt{2m+1}}\sin(p\theta_k),
\]

the rectangular polar identity gives

\begin{equation}\label{eq:supp-leverage-cosine-formula}
 \gamma_{m,j}=\frac1{2m+1}\sum_{k=1}^m
 \frac{1+\cos[(2j-1)\theta_k]}{\sin(\theta_k/2)}.
\end{equation}

Uniformly in \(k\),
\(\csc(\theta_k/2)=2(2m+1)/[\pi(2k-1)]+O(1)\).  Put
\(x_j=(2j-1)\pi/(2m+1)\).  Abel summation and

\[
 \left|\sum_{k=1}^K\cos((2k-1)x)\right|
 \le \min\{K,|\sin x|^{-1}\}
\]

show that the truncated odd-harmonic sum differs from its infinite series by
at most
\[
 \frac{C}{m\sin x_j}=O(1),
\]
uniformly in \(j\).  Hence the identity

\[
 \sum_{k\ge1}\frac{\cos[(2k-1)x]}{2k-1}
 =\frac12\log\cot(x/2),\qquad0<x<\pi,
\]

implies

\[
 \sum_{k=1}^m\frac{\cos[(2k-1)x_j]}{2k-1}
 =\frac12\log\cot(x_j/2)+O(1).
\]

Together with
\(\sum_{k=1}^m(2k-1)^{-1}=\tfrac12\log m+O(1)\) and
\(\log\cot(x_j/2)\le\log_+(Cm/j)\), this proves
\eqref{eq:supp-leverage-profile}, including \(j=1,m\).  Summing the singular
values gives
\eqref{eq:supp-cm-trace}.

For the physical-coordinate diagonal, the exact expansion is

\begin{equation}\label{eq:supp-cm-diagonal-expansion}
 (C_m)_{xx}=\frac1{2m+1}\sum_{k=1}^m
 \frac{1-\cos(2x\theta_k)}{\sin(\theta_k/2)}.
\end{equation}

Let \(\mathsf D_{xk}=\cos(2x\theta_k)\).  Discrete odd-cosine
orthogonality gives

\begin{equation}\label{eq:supp-odd-cosine-gram}
 \mathsf D^{\mathsf T}\mathsf D
 =\frac{2m+1}{4}I-\frac12J.
\end{equation}

Let \(a_k=\csc(\theta_k/2)\), and let \(d_x=(C_m)_{xx}\).  Centering can only
decrease the Euclidean norm, so

\[
 \frac1m\sum_{x=1}^m\left|d_x-\frac1m\sum_{y=1}^md_y\right|^2
 \le
 \frac{a^{\mathsf T}\mathsf D^{\mathsf T}\mathsf D a}
 {m(2m+1)^2}
 \le\frac{\|a\|_2^2}{4m(2m+1)}=O(1),
\]

where the last step uses
\(\sum_k\csc^2(\theta_k/2)=O(m^2)\).  Since the centered mean is
\(\operatorname{tr}C_m/m\), Cauchy--Schwarz now gives
\eqref{eq:supp-cm-flatness}.
\end{proof}

\begin{proposition}[Quantitative Volterra certificate]
\label{prop:supp-quantitative-volterra}
For all sufficiently large \(m\), set

\begin{equation}\label{eq:supp-moving-boundary}
 \delta_m=\frac{3\log\log m}{\log m},
 \qquad h_m=\lceil m^{1-\delta_m}\rceil.
\end{equation}

There is \(K_m\succeq0\) dominating every suffix projector and satisfying

\begin{align}
 \operatorname{tr}K_m
 &\le\frac{m(\log m)^2}{\pi^2}
 \left[1+C\frac{\log\log m}{\log m}\right],
 \label{eq:supp-quantitative-volterra-trace}\\
 \frac1{\operatorname{tr}K_m}\sum_x
 \left|(K_m)_{xx}-\frac{\operatorname{tr}K_m}{m}\right|
 &\le\frac C{\log m}.
 \label{eq:supp-quantitative-volterra-flat}
\end{align}
\end{proposition}

\begin{proof}
Let \(t_m=\max_{j\ge h_m}\gamma_{m,j}\).  Lemma
\ref{lem:supp-odd-cosine} gives
\(t_m\le(1+\delta_m)\log m/\pi+C\), and
\(t_mC_m\) dominates all suffixes beginning at \(j\ge h_m\).
For \(j<h_m\), decompose

\[
 r_{m,j}=w_m+\ell_{m,j},
 \qquad
 w_m=\mathbf1_{\{h_m,\ldots,m\}},
 \quad
 \ell_{m,j}=\mathbf1_{\{j,\ldots,h_m-1\}}.
\]

The cyclic certificate \eqref{eq:supp-cyclic-certificate}, embedded in the
first \(h_m-1\) coordinates, supplies \(K_{h_m}\) with
\(\operatorname{tr}K_{h_m}=O(h_m\log^2h_m)\).  Therefore

\begin{equation}\label{eq:supp-volterra-construction}
 K_m=t_mC_m+2\left(|w_m\rangle\langle w_m|+K_{h_m}\right)
\end{equation}

dominates every suffix.  The chosen boundary scale obeys

\[
 m^{-\delta_m}=O((\log m)^{-3}).
\]

Thus the repair has relative trace cost \(O((\log m)^{-2})\).
Equation~\eqref{eq:supp-cm-trace} proves
\eqref{eq:supp-quantitative-volterra-trace}.  The diagonal estimate
\eqref{eq:supp-cm-flatness} then proves
\eqref{eq:supp-quantitative-volterra-flat}.  Feasibility and
Lemma~\ref{lem:supp-subnormalized-pgm} provide the lower estimate
\(\operatorname{tr}K_m\ge c_0m\log^2m\), for an absolute constant
\(c_0>0\), used in the relative bound.
\end{proof}

\begin{corollary}[Fixed-boundary Volterra certificate]
\label{cor:supp-fixed-delta-volterra}
Fix \(0<\delta<1\) and put
\(h_{m,\delta}=\lceil m^{1-\delta}\rceil\).  The same construction gives
a certificate \(K_{m,\delta}\) dominating every suffix projector and
satisfying

\begin{align}
 \operatorname{tr}K_{m,\delta}
 &\le\left[\frac{1+\delta}{\pi^2}+o_\delta(1)\right]
 m(\log m)^2,
 \label{eq:supp-fixed-delta-trace}\\
 \sum_x\left|(K_{m,\delta})_{xx}
 -\frac{\operatorname{tr}K_{m,\delta}}m\right|
 &=o_\delta(\operatorname{tr}K_{m,\delta}).
 \label{eq:supp-fixed-delta-flatness}
\end{align}

Here \(o_\delta(1)\to0\) as \(m\to\infty\) with \(\delta\) fixed.
\end{corollary}

\begin{proof}
Use \eqref{eq:supp-volterra-construction} with
\(h_m=h_{m,\delta}\).  The leverage bound becomes
\(t_m\le(1+\delta)\log m/\pi+C\).  The repair trace is
\(O(m)+O(m^{1-\delta}\log^2m)=o_\delta(m\log^2m)\).  The odd-cosine
diagonal deviation of \(t_mC_m\) is
\(t_mO(m)=O_\delta(m\log m)=o_\delta(m\log^2m)\), and the repair has
diagonal deviation at most twice its trace.  The lower bound
\(\operatorname{tr}K_{m,\delta}\ge c_0m\log^2m\), for an absolute constant
\(c_0>0\), converts both additive remainders into the stated relative
remainder.
\end{proof}

\begin{lemma}[Singular-safe weighted block and tail extension]
\label{lem:supp-weighted-block-tail}
Let \(Q\succeq0\) on a finite-dimensional Hilbert space and suppose

\[
 Q\succeq|x_a\rangle\langle x_a|,\qquad
 \operatorname{tr}Q\le L,\qquad
 u_a=x_a+\gamma_aw,
 \qquad |\gamma_a|\le1,\quad\|w\|\le1
\]

for every label \(a\).  If \(L>0\), then

\[
 \mathcal K_L=(1+L^{-1/2})Q+(1+\sqrt L)|w\rangle\langle w|
\]

dominates every \(|u_a\rangle\langle u_a|\) and satisfies

\[
 \operatorname{tr}\mathcal K_L\le(\sqrt L+1)^2.
\]

If \(L=0\), then \(Q=0\), every \(x_a=0\), and
\(\mathcal K_0=|w\rangle\langle w|\) is feasible with trace at most one.

For the physical right suffix

\[
 u^+_{a,c}(k)=s(c)c^{k-a}\mathbf1_{\{k\ge a\}}
\]

and a consecutive block \(E=[A,B]\) containing \(a\), the decomposition

\[
 u^+_{a,c}=x^+_{E,a}+c^{B+1-a}w_E^+
\]

is exact, where \(x^+_{E,a}\) is the restriction to \(k\le B\) and

\[
 w_E^+(k)=s(c)c^{k-B-1}\mathbf1_{\{k\ge B+1\}}.
\]

Here \(0\le c^{B+1-a}\le1\) and \(\|w_E^+\|\le1\).  Reflection gives the
left suffix decomposition

\[
 u^-_{b,c}=x^-_{E,b}+c^{b-A+1}w_E^-.
\]
\end{lemma}

\begin{proof}
For arbitrary vectors \(x,y\) and \(\alpha>0\),

\[
 \alpha|x\rangle\langle x|+\alpha^{-1}|y\rangle\langle y|
 -|x\rangle\langle y|-|y\rangle\langle x|
 =
 |\sqrt\alpha\,x-\alpha^{-1/2}y\rangle
 \langle\sqrt\alpha\,x-\alpha^{-1/2}y|\succeq0.
\]

Apply this identity with \(y=\gamma_aw\).  Since
\(|w\rangle\langle w|\succeq|\gamma_aw\rangle
\langle\gamma_aw|\), choosing \(\alpha=L^{-1/2}\) proves feasibility for
\(L>0\), and

\[
 \operatorname{tr}\mathcal K_L
 \le(1+L^{-1/2})L+(1+\sqrt L)=(\sqrt L+1)^2.
\]

If \(L=0\), positivity and trace zero imply \(Q=0\), hence every \(x_a=0\);
the stated direct branch follows.

The suffix identity is checked coordinatewise.  For \(k\ge B+1\), its tail
term is
\(s(c)c^{B+1-a}c^{k-B-1}=s(c)c^{k-a}\), and the coordinates within the block
agree by definition. Moreover,

\[
 \|w_E^+\|^2=(1-c^2)\sum_{j=0}^{n-B-1}c^{2j}\le1.
\]

Reflection about the midpoint of the site set gives the left formula.  At
\(c=0\), the suffix is a coordinate vector and the decomposition is read
directly; at \(c=1\), every physical suffix is zero.
\end{proof}

\begin{proposition}[Weighted local Volterra certificate]
\label{prop:supp-weighted-volterra-certificate}
Let \(K_m^{(0)}\) denote the unweighted certificate supplied by
Proposition~\ref{prop:supp-quantitative-volterra}, and let

\[
 E_{m,c}=\operatorname{diag}(1,c,\ldots,c^{m-1}).
\]

For \(0<c<1\), the operator

\[
 Q_{m,c}=s(c)^2c^{-2(m-1)}
 E_{m,c}K_m^{(0)}E_{m,c}
\]

dominates the rank-one operator associated with every physical suffix
restricted to the block. For all sufficiently large \(m\), it has the explicit
upper budget

\[
 \operatorname{tr}Q_{m,c}\le\Lambda_{m,c}:=
 \frac{s(c)^2c^{-2(m-1)}}{\pi^2}
 m(\log m)^2
 \left[1+C\frac{\log\log m}{\log m}\right].
\]

For any sequence of pairs \((m_j,c_j)\) satisfying

\[
 m_j\to\infty,\qquad 0<c_j<1,\qquad c_j\to1,
 \qquad m_j(1-c_j)\to0,
 \qquad \frac{\log m_j}{\ell(c_j)}\to1,
\]

one has

\[
 \Lambda_{m_j,c_j}=m_jp_1(c_j)[1+o(1)].
\]
More precisely, let
\(\{(m_{n,k},c_{n,k}):k\in K_n\}\) be a triangular array with
\(0<c_{n,k}<1\). If

\[
\begin{aligned}
\inf_{k\in K_n}m_{n,k}&\longrightarrow\infty,
&\sup_{k\in K_n}(1-c_{n,k})&\longrightarrow0,\\
\sup_{k\in K_n}m_{n,k}(1-c_{n,k})&\longrightarrow0,
&\sup_{k\in K_n}
\left|\frac{\log m_{n,k}}{\ell(c_{n,k})}-1\right|&\longrightarrow0,
\end{aligned}
\]

then the conclusion is uniform in \(k\):

\[
\sup_{k\in K_n}
\left|
\frac{\Lambda_{m_{n,k},c_{n,k}}}
{m_{n,k}p_1(c_{n,k})}-1
\right|\longrightarrow0.
\]

At \(c=0\), the identity operator is a direct local certificate; at \(c=1\),
the zero operator is the direct certificate.  Finitely many block sizes not
covered by the quantitative estimate are handled by the finite frame sum,
with its actual trace used as the budget.
\end{proposition}

\begin{proof}
In local coordinates \(j=1,\ldots,m\), let
\(r_{m,j}=\mathbf1_{\{j,\ldots,m\}}\). The physical suffix restricted to the
block is

\[
 x_j=s(c)c^{-(j-1)}E_{m,c}r_{m,j}.
\]

Congruence of
\(K_m^{(0)}\succeq|r_{m,j}\rangle\langle r_{m,j}|\) by \(E_{m,c}\),
together with
\(c^{-2(m-1)}\ge c^{-2(j-1)}\), proves the domination.  Because
\(0\preceq E_{m,c}^2\preceq I\),

\[
 \operatorname{tr}(E_{m,c}K_m^{(0)}E_{m,c})
 =\operatorname{tr}(K_m^{(0)}E_{m,c}^2)
 \le\operatorname{tr}K_m^{(0)}.
\]

This proves the displayed budget. Its ratio to \(mp_1(c)\) is exactly

\[
 \frac{\Lambda_{m,c}}{mp_1(c)}
 =c^{-2(m-1)}
 \left(\frac{\log m}{\ell(c)}\right)^2
 \left[1+C\frac{\log\log m}{\log m}\right].
\]

For the stated sequences this tends to one, since
\(c_j^{-2(m_j-1)}\to1\); the same calculation is uniform for every triangular
family on which the four errors in the proposition tend to zero uniformly.
The endpoint constructions are defined directly and avoid inverse powers of
\(c\).
\end{proof}

\begin{corollary}[Weighted endpoint-cell certificate]
\label{cor:supp-weighted-cell-certificate}
Fix \(n\) and a common overlap \(c\in[0,1]\), and partition the sites into
consecutive endpoint blocks \(E_1,\ldots,E_B\). For each block \(E_p\), let
\(\Lambda_p\) be the applicable local trace budget from
Proposition~\ref{prop:supp-weighted-volterra-certificate}, using its direct
endpoint or finite-frame budget in the exceptional branches.
For each block, extend the local operators of
Proposition~\ref{prop:supp-weighted-volterra-certificate} by zero and apply
Lemma~\ref{lem:supp-weighted-block-tail}.  This gives right and left suffix
certificates \(\mathcal K_p^+\) and \(\mathcal K_p^-\) satisfying

\[
 \operatorname{tr}\mathcal K_p^\pm
 \le T_p:=(\sqrt{\Lambda_p}+1)^2.
\]

Let \(P_\Delta\) restrict ordered site pairs to
\(\Delta_n=\{(u,v):u<v\}\), and define

\[
 \mathcal C_{p,q}
 =P_\Delta(\mathcal K_p^+\otimes\mathcal K_q^-)P_\Delta^*,
 \qquad
 \mathcal K_{\ge2,n}=\sum_{p\le q}\mathcal C_{p,q}.
\]

Then \(\mathcal K_{\ge2,n}\) dominates the physical hull projector of every
interval and

\[
 \operatorname{tr}\mathcal K_{\ge2,n}
 \le\sum_{p\le q}T_pT_q.
\]

The operator \(\mathcal K_{\ge2,n}\) is the unscaled geometric certificate;
the equiprior Bayes dual is
\(\mathcal K_{\ge2,n}/M_n\).
\end{corollary}

\begin{proof}
If \(A\succeq|x\rangle\langle x|\) and
\(B\succeq|y\rangle\langle y|\), with \(A,B\succeq0\), then

\[
 A\otimes B-|x\otimes y\rangle\langle x\otimes y|
 =(A-|x\rangle\langle x|)\otimes B
 +|x\rangle\langle x|\otimes
  (B-|y\rangle\langle y|)\succeq0.
\]

Compressing both sides preserves the domination relation. Every interval has
its endpoints in a unique cell \(p\le q\). The exact hull identity proves
feasibility there, and adding the remaining positive cells cannot weaken it.
The diagonal cells
are also compressed to \(u<v\), never replaced by square cells.  Compression
decreases the trace of a positive operator, giving the trace bound and the
stated normalization.
\end{proof}

\subsection{Triangular compression and the sector bound}

Recall that \(m=n-1\) and

\[
 \Lambda_m=\{(x,y)\in\{1,\ldots,m\}^2:x+y\le m+1\}.
\]

The excitation pair \(s<t\) maps to \((s,n+1-t)\), while an interval
\([a,b]\), \(a<b\), maps to \((a,n+1-b)\).  The corresponding unweighted
containment vector is

\begin{equation}\label{eq:supp-triangular-vector}
 z_{u,v}=P_{\Lambda_m}(r_{m,u}\otimes r_{m,v}).
\end{equation}

\begin{lemma}[Flat-diagonal triangular trace]
\label{lem:supp-triangular-trace}
If \(K_m\succeq0\), \(T_m=\operatorname{tr}K_m\), and
\(\sum_x|(K_m)_{xx}-T_m/m|\le\epsilon_mT_m\), then

\begin{equation}\label{eq:supp-triangular-trace}
 \operatorname{tr}\!\left[
 P_{\Lambda_m}(K_m\otimes K_m)P_{\Lambda_m}\right]
 =\frac12T_m^2[1+O(\epsilon_m+m^{-1})].
\end{equation}
\end{lemma}

\begin{proof}
For \(k_x=(K_m)_{xx}\) and \(\bar k=T_m/m\),
\(\sum_{x,y}|k_xk_y-\bar k^2|\le2\epsilon_mT_m^2\).  Sum over
\(\Lambda_m\) and use \(|\Lambda_m|/m^2=1/2+O(m^{-1})\).
\end{proof}

With \(K_m\) from Proposition~\ref{prop:supp-quantitative-volterra}, define

\begin{equation}\label{eq:supp-triangular-dual}
 \mathcal L_n=P_{\Lambda_m}(K_m\otimes K_m)P_{\Lambda_m}.
\end{equation}

Tensor order and compression give \(\mathcal L_n\succeq|z_{u,v}\rangle\langle
z_{u,v}|\), and

\begin{equation}\label{eq:supp-triangular-dual-trace}
 \operatorname{tr}\mathcal L_n
 \le\frac{n^2(\log n)^4}{2\pi^4}
 \left[1+C\frac{\log\log n}{\log n}\right].
\end{equation}

The exact physical hull vector is
\(q_n^2c_n^{|I|}W_nz_I\).  Hence

\begin{equation}\label{eq:supp-physical-hull-dual}
 \widehat{\mathcal L}_n=q_n^4W_n\mathcal L_nW_n
\end{equation}

is feasible for every depth-\(\ge2\) sector label, because
\(c_n^{|I|}\le1\).  Moreover,
\(\operatorname{tr}\widehat{\mathcal L}_n
\le q_n^4\rho_n^2\operatorname{tr}\mathcal L_n\).

\begin{theorem}[Quantitative sectorwise upper bound]
\label{thm:supp-sector-upper}
Uniformly for \(0\le\tau_n\le T\),

\begin{align}
 M_nP_{\rm opt}^{(0)}&=c_n^2,\label{eq:supp-vacuum-optimum}\\
 M_nP_{\rm opt}^{(1)}
 &\le\frac{4\tau_n}{\pi^2}
 \left[1+O_T((\log n)^{-1})\right],
 \label{eq:supp-one-optimum}\\
 M_nP_{\rm opt}^{(\ge2)}
 &\le\frac{2\tau_n^2}{\pi^4}
 \left[1+O_T\!\left(\frac{\log\log n}{\log n}\right)\right].
 \label{eq:supp-two-optimum}
\end{align}

Consequently,

\begin{equation}\label{eq:supp-full-optimum-upper}
 \sqrt{M_nP_{\rm opt}(G_n)}
 \le A(\tau_n)+C_T\frac{\log\log n}{\log n}.
\end{equation}
\end{theorem}

\begin{proof}
The vacuum vectors are collinear, proving \eqref{eq:supp-vacuum-optimum}.
Equation \eqref{eq:supp-cyclic-certificate}, multiplied by \(s_n^2\), and
\[
 s_n^2n(\log n)^2
 =2\tau_n\left[1+O_T\!\left((n(\log n)^2)^{-1}\right)\right],
\]
prove
\eqref{eq:supp-one-optimum}.  Equations
\eqref{eq:supp-triangular-dual-trace}--
\eqref{eq:supp-physical-hull-dual}, together with
\[
 q_n^2n(\log n)^2
 =2\tau_n\left[1+O_T\!\left((n(\log n)^2)^{-1}\right)\right],
\]
and
\(\rho_n=1+O_T((\log n)^{-2})\), prove
\eqref{eq:supp-two-optimum}.  Finally apply the orthogonal-sector Minkowski
inequality from the main text to the sectors \(0,1,\ge2\).
\end{proof}

\section{Uniform matching and hull dominance}
\label{supp:matching}

\subsection{Inner matching}

\begin{theorem}[Inner matching below the continuum scale]
\label{thm:supp-inner-matching}
Suppose

\begin{equation}\label{eq:supp-inner-assumptions}
 \tau_n=n(1-c_n)(\log n)^2\longrightarrow\infty,
 \qquad
 \lambda_n=n(1-c_n)=\frac{\tau_n}{(\log n)^2}\longrightarrow0.
\end{equation}

Then, for \(X\in\{\mathrm{tr},\mathrm{SRM},\mathrm{opt}\}\),

\begin{equation}\label{eq:supp-inner-conclusion}
 M_nP_X(G_n)\sim\frac{2\tau_n^2}{\pi^4}
 \sim A(\tau_n)^2.
\end{equation}
\end{theorem}

\begin{proof}
Equations \eqref{eq:supp-inner-assumptions} imply
\(c_n^n\to1\) and \(\rho_n=c_n^{-(n-2)}\to1\).  Schatten ideal
inequalities applied to the exact hull factorization give

\begin{equation}\label{eq:supp-inner-schatten-sandwich}
 c_n^n\|Z_n\|_{S_1}
 \le\|D_nZ_nW_n\|_{S_1}
 \le\rho_n\|Z_n\|_{S_1}.
\end{equation}

For example, the lower inequality follows by writing
\(Z_n=D_n^{-1}(D_nZ_nW_n)W_n^{-1}\) and using
\(\|D_n^{-1}\|_{\rm op}=c_n^{-n}\) and
\(\|W_n^{-1}\|_{\rm op}\le1\).  Proposition
\ref{prop:supp-b2-trace} and

\[
 q_n^2n(\log n)^2=2\tau_n[1+o(1)]
\]

therefore give the full physical \(\ge2\)-sector trace amplitude

\begin{equation}\label{eq:supp-inner-sector-amplitude}
 \frac{\|A_{\ge2,n}\|_{S_1}}{\sqrt{M_n}}
 \sim\frac{\sqrt2\,\tau_n}{\pi^2}.
\end{equation}

For the reverse sector bound, use
Corollary~\ref{cor:supp-fixed-delta-volterra} with any fixed
\(0<\delta<1\), rather than the moving choice in
\eqref{eq:supp-moving-boundary}.  Its triangular tensor trace is at most

\[
 \left[\frac{(1+\delta)^2}{2\pi^4}+o_\delta(1)\right]
 n^2(\log n)^4,
\]

where \(o_\delta(1)\to0\) as \(n\to\infty\) with \(\delta\) fixed.
Congruence by
\(q_n^2W_n\) gives the explicit physical bound

\[
 M_nP_{\rm opt}^{(\ge2)}
 \le q_n^4\rho_n^2
 \left[\frac{(1+\delta)^2}{2\pi^4}+o_\delta(1)\right]
 n^2(\log n)^4.
\]

Moreover,

\[
 \frac{q_n^4n^2(\log n)^4}{4\tau_n^2}\longrightarrow1,
 \qquad \rho_n\longrightarrow1.
\]

Therefore

\[
 \limsup_n
 \frac{M_nP_{\rm opt}^{(\ge2)}}{2\tau_n^2/\pi^4}
 \le(1+\delta)^2.
\]

First take \(n\to\infty\), then \(\delta\downarrow0\), and combine with
the subnormalized trace lower bound to obtain equality at leading order.
The vacuum sector is \(O(1)\), while the cyclic certificate makes the
one-excitation sector amplitude \(O(\sqrt{\tau_n})\).  They are negligible
relative to \(\tau_n\).  Orthogonal-sector Minkowski and
\(P_{\rm tr}\le P_{\rm SRM}\le P_{\rm opt}\) prove
\eqref{eq:supp-inner-conclusion}.
\end{proof}

To compare with the fixed-overlap expression, the complete elliptic-integral
form of the one-dimensional functional gives, for
\(\varepsilon=1-c\downarrow0\),

\begin{align}
 p_1(c)
 &=\frac{2\varepsilon}{\pi^2}
 \log^2\!\frac8\varepsilon\,[1+o(1)],
 \label{eq:supp-p1-endpoint}\\
 p_1(c)^2
 &=\frac{4\varepsilon^2}{\pi^4}
 \log^4\!\frac8\varepsilon\,[1+o(1)].
 \label{eq:supp-p1-square-endpoint}
\end{align}

Under \eqref{eq:supp-inner-assumptions}, one has
\(0<\varepsilon_n=1-c_n=\tau_n/[n(\log n)^2]\to0\), and hence

\begin{equation}\label{eq:supp-outer-formula}
 M_np_1(c_n)^2
 \sim\frac{2\tau_n^2}{\pi^4}
 \left[\frac{\log(8n(\log n)^2/\tau_n)}{\log n}\right]^4.
\end{equation}

More generally, for sequences with \(0<1-c_n\to0\), the formulas match the
inner coefficient precisely when
\(\log[n(1-c_n)]=o(\log n)\). This condition holds automatically under
\eqref{eq:supp-inner-assumptions}. Equation~\eqref{eq:supp-outer-formula}
compares the two asymptotic formulas directly through the endpoint expansion.

\subsection{Continuum and moving outer regimes}

When \(\lambda_n=n(1-c_n)\) stays away from zero, the physical length and hull
weights are no longer close to the identity. We therefore retain them within
each block through the exact hull coordinates, local block and tail
certificates, and matched rectangular pinching instead of using
\eqref{eq:supp-inner-schatten-sandwich}.

\begin{lemma}[Matched rectangular pinching]
\label{lem:supp-weighted-rectangular-pinching}
Let \(A:\mathbb C^d\to\mathbb C^r\).  Let
\(\{P_j\}_{j=1}^J\) and \(\{Q_j\}_{j=1}^J\) be pairwise orthogonal
families of row and column projections.  Then

\[
 \left\|\sum_{j=1}^J P_jAQ_j\right\|_{S_1}
 \le\|A\|_{S_1}.
\]
\end{lemma}

\begin{proof}
Use the Hermitian dilation

\[
 \mathcal H(A)=
 \begin{pmatrix}0&A\\A^*&0\end{pmatrix},
 \qquad \|\mathcal H(A)\|_{S_1}=2\|A\|_{S_1}.
\]

The projections \(R_j=P_j\oplus Q_j\), together with the complementary
projection, form an orthogonal resolution of the dilation space.  Ordinary
block pinching is contractive in the Schatten \(1\)-norm, and compression to
\(\sum_jR_j\) leaves the dilation of \(\sum_jP_jAQ_j\).  Dividing the
resulting inequality by two proves the claim.  This argument is specific to
matched row and column blocks; arbitrary entry deletion need not be
contractive.
\end{proof}

\begin{theorem}[Compact continuum interface]
\label{thm:supp-continuum-interface}
Let \(K=[\lambda_-,\lambda_+]\Subset(0,\infty)\) and
\(c_{n,\lambda}=1-\lambda/n\).  For
\(X\in\{\mathrm{tr},\mathrm{SRM},\mathrm{opt}\}\),

\[
 \sup_{\lambda\in K}
 \left|
 \frac{P_X(G_n(c_{n,\lambda}))}
      {p_1(c_{n,\lambda})^2}-1
 \right|\longrightarrow0.
\]

Equivalently,

\[
 \sup_{\lambda\in K}
 \left|
 \frac{M_nP_X(G_n(c_{n,\lambda}))}{(\log n)^4}
 -\frac{2\lambda^2}{\pi^4}
 \right|\longrightarrow0.
\]
\end{theorem}

\begin{proof}
Write \(c=1-\lambda/n\), \(s=s(c)\), and
\(B=\lceil\ell(c)\rceil\).  Partition the sites into \(B\) balanced
consecutive blocks of sizes \(m_p\in\{\lfloor n/B\rfloor,
\lceil n/B\rceil\}\).  Lemma~\ref{lem:supp-elliptic-endpoint} gives, uniformly
for \(\lambda\in K\),

\[
 \ell(1-\lambda/n)=\log n+\log(8/\lambda)+o_K(1).
\]

Consequently,

\[
 B\to\infty,\quad B/n\to0,\quad
 \min_p m_p\to\infty,\quad
 \max_p m_p(1-c)\to0,\quad
 \max_p\left|\frac{\log m_p}{\ell(c)}-1\right|\to0.
\]

The finite-dimensional Volterra law
\eqref{eq:supp-volterra-trace} and the exact similarity

\[
 V_m(c)=E_{m,c}^{-1}U_mE_{m,c}
\]

therefore imply, uniformly over all blocks,

\[
 s\|V_{m_p}(c)\|_{S_1}
 =m_p\sqrt{p_1(c)}[1+o_K(1)].
\]

Indeed, the two ideal-norm comparisons differ from one by at most
\(\exp\{O(m_p(1-c))\}\), and
\(\sqrt{p_1(c)}=s\ell(c)/\pi\).

For each \(p<q\), retain rows with endpoints in
\(E_p\times E_q\) and strict hull columns in the same rectangle.  By
\eqref{eq:supp-weighted-hull-identity}, the retained matrix is exactly

\[
 s^2V_{m_p}(c)\otimes V_{m_q}(c)^{\mathsf T}.
\]

Lemma~\ref{lem:supp-weighted-rectangular-pinching} and tensor
multiplicativity give

\[
 \|A_{\ge2,n}(c)\|_{S_1}
 \ge p_1(c)[1-o_K(1)]\sum_{p<q}m_pm_q.
\]

The restriction to block pairs \(p<q\) gives the exact identity

\[
 \sum_{p<q}m_pm_q
 =\frac12\left(n^2-\sum_pm_p^2\right)
 =M_n[1-o_K(1)],
\]

because \(\sum_pm_p^2\le n(n/B+1)\).  Thus

\[
 \|A_{\ge2,n}(c)\|_{S_1}
 \ge M_np_1(c)[1-o_K(1)].
\]

The full Gram matrix dominates
\(A_{\ge2,n}A_{\ge2,n}^*\).  Operator monotonicity of the square root,
followed by the full external normalization \(1/M_n\), yields

\[
 P_{\rm tr}(G_n(c))
 \ge\left(\frac{\|A_{\ge2,n}(c)\|_{S_1}}{M_n}\right)^2
 \ge p_1(c)^2[1-o_K(1)].
\]

For the matching upper bound, Proposition
\ref{prop:supp-weighted-volterra-certificate} gives local budgets

\[
 \Lambda_{p,n}=m_pp_1(c)[1+o_K(1)].
\]

Their minimum diverges: uniformly on \(K\),
\(m_pp_1(c)\asymp_K\ell(c)\to\infty\).  Lemma
\ref{lem:supp-weighted-block-tail} therefore gives full one-sided suffix
certificates with

\[
 T_{p,n}:=(\sqrt{\Lambda_{p,n}}+1)^2
 =m_pp_1(c)[1+o_K(1)].
\]

Corollary~\ref{cor:supp-weighted-cell-certificate}, including all diagonal
cells, supplies an unscaled geometric hull certificate satisfying

\[
 \operatorname{tr}\mathcal K_{\ge2,n}
 \le\sum_{p\le q}T_{p,n}T_{q,n}
 \le M_np_1(c)^2[1+o_K(1)].
\]

Dividing this geometric certificate by \(M_n\) yields
\(P_{\rm opt}^{(\ge2)}\le p_1(c)^2[1+o_K(1)]\).

The vacuum sector is exact:
\(P_{\rm opt}^{(0)}=c^2/M_n\).  The cyclic certificate
\eqref{eq:supp-cyclic-certificate} gives

\[
 \kappa_{1,n}\le C s^2n(\log n)^2.
\]

The compact scalar normalization then yields, uniformly on \(K\),

\[
 \frac{\sqrt{P_{\rm opt}^{(0)}}}{p_1(c)}
 =O_K((\log n)^{-2}),\qquad
 \frac{\sqrt{P_{\rm opt}^{(1)}}}{p_1(c)}
 =O_K((\log n)^{-1}).
\]

Orthogonal-sector Minkowski, the hull certificate, and
\(P_{\rm tr}\le P_{\rm SRM}\le P_{\rm opt}\) squeeze all three relative
ratios to one.

Finally, the exact scalar identity

\[
 \frac{M_np_1(1-\lambda/n)^2}{(\log n)^4}
 =\frac{\lambda^2}{2\pi^4}
 \left(1+\frac1n\right)
 \left(2-\frac{\lambda}{n}\right)^2
 \left(\frac{\ell(1-\lambda/n)}{\log n}\right)^4
\]

converges uniformly to \(2\lambda^2/\pi^4\), and is uniformly bounded
away from zero on \(K\). This proves the scaled formulation.
\end{proof}

The table records the normalization conventions used here and below.

\begin{center}
\small
\begin{tabularx}{\linewidth}{@{}lXl@{}}
\toprule
Object & Constraint or definition & Prior status\\
\midrule
\(A_{\ge2,n}\) & Exact physical strict-hull state matrix; all \(M_n\)
rows, including zero singleton rows & no prior\\
\(\mathcal K_{\ge2,n}\) &
\(\mathcal K_{\ge2,n}\succeq
 |\phi_I^{(\ge2)}\rangle\langle\phi_I^{(\ge2)}|\) for every
physical interval & geometric, unscaled\\
\(\mathcal K_{\ge2,n}/M_n\) & Feasible Bayes dual for the hull sector &
external \(1/M_n\)\\
\(P_{\rm opt}^{(k)}\) & \(\kappa_{k,n}/M_n\), for
\(k\in\{0,1,\ge2\}\) & same \(1/M_n\) in every sector\\
\(P_{\rm tr}(G_n)\) &
\((\operatorname{tr}\sqrt{G_n}/M_n)^2\) for the full physical Gram &
full ensemble\\
\bottomrule
\end{tabularx}
\end{center}

\begin{lemma}[Adaptive macroblock parameters]
\label{lem:supp-adaptive-block-parameters}
Suppose \(c_n\to1\), \(c_n<1\) eventually, and
\(h_n=np_1(c_n)\to\infty\).  Set

\[
 B_n=\left\lceil\sqrt{h_n}+\lambda_n\sqrt{\ell_n}\right\rceil,
 \qquad m_n=\left\lfloor\frac n{B_n}\right\rfloor.
\]

Then

\[
 B_n\to\infty,\qquad
 \frac{B_n}{h_n}\to0,\qquad
 \frac{B_n}{n}\to0,\qquad
 \frac{\lambda_n}{B_n}\to0,
\]

and a balanced \(B_n\)-block partition exists eventually.  Its minimum
block size satisfies

\[
 m_n(1-c_n)\to0,\qquad
 m_np_1(c_n)\to\infty,\qquad
 \frac{\log m_n}{\ell_n}\to1.
\]
No monotonicity of \(c_n\) or \(\lambda_n\) is required.
\end{lemma}

\begin{proof}
Cauchy--Schwarz and the Poisson-kernel integral give
\(J(c)^2\le(1-c^2)^{-1}\), hence \(p_1(c)\le1\) and \(h_n\le n\).
The exact identity

\[
 h_n=\frac{\lambda_n(1+c_n)}{\pi^2}\ell_n^2
\]

and the definition of \(B_n\) give

\[
 \frac{B_n}{h_n}
 \le \frac{1}{\sqrt{h_n}}
 +\frac{\pi^2}{(1+c_n)\ell_n^{3/2}}
 +\frac{1}{h_n}\longrightarrow0.
\]

Also \(B_n\ge\sqrt{h_n}\to\infty\),
\(\lambda_n/B_n\le\ell_n^{-1/2}\to0\), and
\(B_n/n\le B_n/h_n\to0\).  Thus \(B_n\le n\) eventually and

\[
 m_n=\frac n{B_n}[1+O(B_n/n)].
\]

Multiplication by \(1-c_n=\lambda_n/n\) and by
\(p_1(c_n)=h_n/n\) proves the first two block limits.

It remains to match logarithms. Lemma~\ref{lem:supp-elliptic-endpoint} gives

\[
 \ell_n=\log n-\log\lambda_n+\log8+o(1).
\]

On indices with \(\lambda_n\le1\), the identity for \(h_n\) and
\(h_n\to\infty\) imply
\(-\log\lambda_n\le2\log\ell_n\) eventually.  Hence
\(\ell_n=\log n+O(\log\ell_n)\), while the block definition gives
\(\log B_n=O(\log\ell_n)\).  On indices with \(\lambda_n\ge1\),

\[
 \lambda_n\sqrt{\ell_n}\le B_n
 \le C\lambda_n\ell_n
\]

eventually, so
\(\log B_n=\log\lambda_n+O(\log\ell_n)\).  In both cases,

\[
 \log(n/B_n)=\ell_n+O(\log\ell_n).
\]

The floor changes this by \(o(1)\), and \(\ell_n\to\infty\), proving
\(\log m_n/\ell_n\to1\). Together, the two cases cover every index, even when
the sequence oscillates.
\end{proof}

\begin{lemma}[Adaptive strict-hull squeeze]
\label{lem:supp-adaptive-hull-squeeze}
Suppose \(c_n\to1\), \(c_n<1\) eventually, and
\(\lambda_n=n(1-c_n)\to\infty\).  The identity
\(h_n=\lambda_n(1+c_n)\ell_n^2/\pi^2\) gives \(h_n\to\infty\).  Use the
balanced partition with \(B_n\) blocks and minimum block size \(m_n\) from
Lemma~\ref{lem:supp-adaptive-block-parameters}, and define

\[
 \varepsilon_n^{(\ge2)}=
 \frac{\lambda_n}{B_n}
 +\frac{\log\log m_n}{\log m_n}
 +\left|\frac{\log m_n}{\ell_n}-1\right|
 +\sqrt{\frac{B_n}{h_n}}
 +\frac1{B_n}.
\]

There is an absolute constant \(C>0\) such that, for all sufficiently large
\(n\), the exact strict-hull state matrix and the certificate constructed in
Corollary~\ref{cor:supp-weighted-cell-certificate} satisfy

\begin{align}
 \|A_{\ge2,n}(c_n)\|_{S_1}
 &\ge M_np_1(c_n)\bigl(1-E_n^{(\ge2)}\bigr),
 \label{eq:supp-adaptive-hull-lower}\\
 \operatorname{tr}\mathcal K_{\ge2,n}
 &\le M_np_1(c_n)^2\bigl(1+E_n^{(\ge2)}\bigr),
 \label{eq:supp-adaptive-hull-upper}
\end{align}

where \(E_n^{(\ge2)}=C\varepsilon_n^{(\ge2)}\to0\).
\end{lemma}

\begin{proof}
For each pair of distinct blocks, the exact weighted-hull identity
\eqref{eq:supp-weighted-hull-identity} and
Lemma~\ref{lem:supp-weighted-rectangular-pinching} retain the tensor product
of the two local weighted Volterra factors.  The similarity comparison costs
\(O(m_n(1-c_n))=O(\lambda_n/B_n)\); the finite Volterra trace law costs
\(O(1/\log m_n)\); and replacing \(\log m_n\) by \(\ell_n\) costs
\(O(|\log m_n/\ell_n-1|)\).  Finally,

\[
 \sum_{p<q}m_{p,n}m_{q,n}
 =\frac12\left(n^2-\sum_pm_{p,n}^2\right)
 =M_n\bigl[1-O(B_n^{-1})\bigr].
\]

These estimates prove \eqref{eq:supp-adaptive-hull-lower}, after enlarging a
single absolute constant \(C\); the term \(1/\log m_n\) is absorbed by
\(\log\log m_n/\log m_n\) for large \(n\).

For the upper bound, Proposition
\ref{prop:supp-weighted-volterra-certificate} gives, uniformly for the two
possible block sizes,

\[
 \frac{\Lambda_{p,n}}{m_{p,n}p_1(c_n)}
 =1+O\!\left(
 \frac{\lambda_n}{B_n}
 +\frac{\log\log m_n}{\log m_n}
 +\left|\frac{\log m_n}{\ell_n}-1\right|
 \right).
\]

Lemma~\ref{lem:supp-weighted-block-tail} changes the local budget to
\(T_{p,n}=(\sqrt{\Lambda_{p,n}}+1)^2\).  Since
\(\min_p\Lambda_{p,n}\asymp h_n/B_n\to\infty\), its relative cost is
\(O(\sqrt{B_n/h_n})\).  Summing the endpoint-cell certificates of
Corollary~\ref{cor:supp-weighted-cell-certificate} over \(p\le q\) adds only
\(O(B_n^{-1})\), because

\[
 \sum_{p\le q}m_{p,n}m_{q,n}
 =\frac12\left(n^2+\sum_pm_{p,n}^2\right)
 =M_n\bigl[1+O(B_n^{-1})\bigr].
\]

This proves \eqref{eq:supp-adaptive-hull-upper}.  Every term in
\(\varepsilon_n^{(\ge2)}\) tends to zero by
Lemma~\ref{lem:supp-adaptive-block-parameters} and its logarithmic estimate,
so \(E_n^{(\ge2)}\to0\).
\end{proof}

\begin{theorem}[Moving outer regime]
\label{thm:supp-moving-outer}
If \(c_n\to1\), \(c_n<1\) eventually, and
\(\lambda_n=n(1-c_n)\to\infty\), then, for
\(X\in\{\mathrm{tr},\mathrm{SRM},\mathrm{opt}\}\),

\[
 \frac{P_X(G_n(c_n))}{p_1(c_n)^2}\longrightarrow1.
\]
\end{theorem}

\begin{proof}
The identity
\(h_n=\lambda_n(1+c_n)\ell_n^2/\pi^2\) gives \(h_n\to\infty\).
Lemma~\ref{lem:supp-adaptive-hull-squeeze} therefore
supplies both strict-hull bounds.  In particular,
\(P_{\rm tr}(G_n)\ge p_1(c_n)^2[1-o(1)]\), while the Bayes-normalized hull
certificate gives
\(P_{\rm opt}^{(\ge2)}\le p_1(c_n)^2[1+o(1)]\).

The vacuum relative amplitude is \(O(h_n^{-1})\).  For one excitation,
\eqref{eq:supp-cyclic-certificate} gives

\[
 \frac{\sqrt{P_{\rm opt}^{(1)}}}{p_1(c_n)}
 =O\!\left(
 \frac{\log n}{\sqrt{\lambda_n}\,\ell_n^2}
 \right)=o(1).
\]

For the last implication, put
\(q_n^{\rm sc}=n/\lambda_n=(1-c_n)^{-1}\to\infty\).
Then \(\ell_n=\log q_n^{\rm sc}+O(1)\), while
\(\log n=\log\lambda_n+\log q_n^{\rm sc}\).  Both
\(\log\lambda_n/\sqrt{\lambda_n}\) and
\((\sqrt{\lambda_n}\log q_n^{\rm sc})^{-1}\) tend to zero.
Define the remaining sector errors by

\[
 E_n^{(1)}=\frac{\log n}{\sqrt{\lambda_n}\ell_n^2},
 \qquad E_n^{(0)}=\frac1{h_n},
 \qquad \mathcal E_n=E_n^{(\ge2)}+E_n^{(1)}+E_n^{(0)}.
\]

The preceding estimates show that \(\mathcal E_n\to0\).  Sector Minkowski and
\(P_{\rm tr}\le P_{\rm SRM}\le P_{\rm opt}\) now complete the squeeze.  Thus
\(\mathcal E_n\) is also a sufficient control, up to an absolute multiplicative
constant, for the relative amplitude error.
\end{proof}

\subsection{Unified hull-dominance criterion}

\begin{theorem}[Unified hull-dominance criterion]
\label{thm:supp-hull-dominance}
Suppose

\[
 c_n\to1,\qquad c_n<1\ \text{eventually},\qquad
 h_n=np_1(c_n)\to\infty.
\]

Then, without a monotonicity assumption and simultaneously for
\(X\in\{\mathrm{tr},\mathrm{SRM},\mathrm{opt}\}\),

\[
 \frac{P_X(G_n(c_n))}{p_1(c_n)^2}\longrightarrow1.
\]
\end{theorem}

\begin{proof}
Fix an arbitrary subsequence and pass to a further subsequence on which
\(\lambda_n=n(1-c_n)\) converges in the compactified half-line
\([0,\infty]\).

If \(\lambda_n\to0\), then the exact identity

\[
 \lambda_n\ell_n^2=\frac{\pi^2h_n}{1+c_n}\longrightarrow\infty
\]

implies \(-\log\lambda_n\le2\log\ell_n\) eventually.  Bootstrapping this
bound in
\(\ell_n=\log n-\log\lambda_n+\log8+o(1)\) gives
\(\ell_n/\log n\to1\).  Therefore

\[
 \tau_n=\lambda_n(\log n)^2
 =\frac{\pi^2h_n}
 {(1+c_n)(\ell_n/\log n)^2}\longrightarrow\infty.
\]

The hypotheses of Theorem~\ref{thm:supp-inner-matching} now hold.  Its
normalization agrees with the present one because

\[
 \frac{M_np_1(c_n)^2}{2\tau_n^2/\pi^4}
 =\frac{M_n}{n^2}\frac{(1+c_n)^2}{2}
  \left(\frac{\ell_n}{\log n}\right)^4
 \longrightarrow1.
\]

If \(\lambda_n\to\lambda_*\in(0,\infty)\), choose one compact
\(K\Subset(0,\infty)\) containing the values eventually and evaluate the
uniform bound in Theorem~\ref{thm:supp-continuum-interface} at the moving
point \(\lambda_n\).  If \(\lambda_n\to\infty\), apply
Theorem~\ref{thm:supp-moving-outer}.

To apply the inner and outer theorems along a sparse subsequence, leave the
selected values unchanged.  At all sufficiently large omitted dimensions, set
\(\lambda_n=1/\log n\) in the inner case or \(\lambda_n=\sqrt n\) in the outer
case, and define

\[
 c_n=1-\frac{\lambda_n}{n}.
\]

Assign arbitrary legal values on the finite remaining prefix. In both
constructions \(0\le c_n<1\) eventually. The extension preserves the selected
Gram matrices and satisfies the corresponding regime assumptions.
The continuum case requires no extension because
Theorem~\ref{thm:supp-continuum-interface} is uniform on compact parameter
intervals.

Thus every subsequence has a further subsequence on which the displayed
ratio tends to one. The standard subsequence criterion proves the full limit.
\end{proof}

The table lists the regimes, their proof inputs, and their common limits.

\begin{center}
\small
\begin{tabularx}{\linewidth}{@{}l
 >{\raggedright\arraybackslash}X
 >{\raggedright\arraybackslash}X
 >{\raggedright\arraybackslash}X@{}}
\toprule
Regime & Scale assumptions & Proof input & Common conclusion\\
\midrule
Ultracritical &
\(\tau_n=n(1-c_n)(\log n)^2\) in a fixed compact window &
global trace theorem plus sector duals &
\(M_nP_X=A(\tau_n)^2+o(1)\)\\
Inner hull dominance &
\(\lambda_n\to0,\ h_n\to\infty\) &
inner bridge and Theorem~\ref{thm:supp-inner-matching} &
\(P_X\sim p_1(c_n)^2\)\\
Continuum &
\(\lambda_n\) in a fixed compact subset of \((0,\infty)\) &
weighted cells and a squeeze uniform on compact parameter intervals &
\(P_X\sim p_1(c_n)^2\), with coefficient
\(2\lambda_n^2/\pi^4\) after \(M_n/(\log n)^4\) scaling\\
Moving outer &
\(c_n\to1,\ c_n<1\) eventually, \(\lambda_n\to\infty\) &
adaptive blocks and weighted block and tail certificates &
\(P_X\sim p_1(c_n)^2\)\\
Unified &
\(c_n\to1,\ c_n<1\) eventually, \(h_n\to\infty\), with arbitrary
oscillation &
compactified subsequence closure &
\(P_X\sim p_1(c_n)^2\)\\
\bottomrule
\end{tabularx}
\end{center}

Every row uses the physical Gram matrix
\eqref{eq:supp-physical-gram}, the same \(M_n\), and the same external uniform
prior. On the canonical fixed-\(\tau\) scale, \(h_n\) remains bounded; the
hull-dominance theorem instead assumes \(h_n\to\infty\).

\section{Independent labelwise SRM refinement}
\label{supp:direct-srm}

The result below is an independent refinement that resolves the SRM
contribution label by label.  It is not used in the global trace and dual
squeeze that proves the headline optimum and SRM laws.

\begin{theorem}[Physical local law on compact parameter intervals]
\label{thm:supp-full-local-law}
For every \(T<\infty\), with
\(c_n(t)=1-t/[n(\log n)^2]\),

\begin{equation}\label{eq:supp-full-local-law}
 \sup_{0\le t\le T}\frac1{M_n}\sum_{I\in\mathcal I_n}
 \left|\sqrt{M_n}(\sqrt{G_n(c_n(t))})_{I,I}-A(t)\right|^2
 \longrightarrow0.
\end{equation}
\end{theorem}

The proof combines the retained orthogonality, incidence-trace, and odd-cosine
estimates established earlier.

\subsection{One- and two-excitation local laws}

We use a finite-codimension estimate when the three excitation strata are
made row-orthogonal. Let
\((U_m)_{ij}=\mathbf1_{\{i\le j\}}\).

\begin{lemma}[Finite-codimension Volterra stability]
\label{lem:supp-finite-codimension}
Let \(Q_m\) be an orthogonal projection with
\(\operatorname{rank}(I-Q_m)\le r\), where \(r\) is fixed, and put

\[
 h_m^{(Q)}(j)=
 \left[(U_mQ_mU_m^*)^{1/2}\right]_{jj}.
\]

Uniformly over all such projections,

\begin{equation}\label{eq:supp-finite-codimension-local}
 \frac1m\sum_{j=1}^m
 \left(\frac{h_m^{(Q)}(j)}{\log m}-\frac1\pi\right)^2
 \longrightarrow0,
 \qquad
 \max_jh_m^{(Q)}(j)\le\frac2\pi\log m+O(1).
\end{equation}
\end{lemma}

\begin{proof}
Write \(h_m=h_m^{(I)}\).  Operator monotonicity gives
\(0\le h_m^{(Q)}(j)\le h_m(j)\), while

\begin{align*}
 0&\le \|U_m\|_{S_1}-\|U_mQ_m\|_{S_1}\\
 &\le \|U_m(I-Q_m)\|_{S_1}
 \le \sqrt r\,\|U_m\|_{S_2}=O_r(m).
\end{align*}

Thus the nonnegative deficit
\((h_m-h_m^{(Q)})/\log m\) has empirical mean
\(O_r((\log m)^{-1})\).  The odd-cosine expansion in
Lemma~\ref{lem:supp-odd-cosine} gives

\[
 \frac1m\sum_j
 \left(\frac{h_m(j)}{\log m}-\frac1\pi\right)^2\to0,
 \qquad
 \max_jh_m(j)\le\frac2\pi\log m+O(1).
\]

The normalized deficit is therefore uniformly bounded; its empirical
second moment is bounded by a constant times its empirical first moment.
This proves both assertions.
\end{proof}

Let
\(H_{1,n}=(B_{1,n}B_{1,n}^*)^{1/2}\) and
\(h_{1,n}(I)=(H_{1,n})_{I,I}\).

\begin{lemma}[One-excitation empirical law]
\label{lem:supp-b1-local}
With \(N=n+1\),

\begin{equation}\label{eq:supp-b1-local}
 \frac1{M_n}\sum_{I\in\mathcal I_n}
 \left[\frac{\sqrt N}{\log N}h_{1,n}(I)-\frac2\pi\right]^2
 \longrightarrow0.
\end{equation}

Moreover,

\begin{equation}\label{eq:supp-b1-local-maximum}
 \max_{I\in\mathcal I_n}h_{1,n}(I)
 \le C\frac{\log N}{\sqrt N}.
\end{equation}
\end{lemma}

\begin{proof}
The Dirichlet eigenbasis in \eqref{eq:supp-b1-normal} gives, for
\(I=[a,b]\) and \(L=b-a+1\),

\begin{equation}\label{eq:supp-b1-exact-diagonal}
 h_{1,n}([a,b])=\frac4{N^{3/2}}\sum_{k=1}^{N-1}
 \frac{\sin^2(L\theta_k/2)\sin^2[(a+b)\theta_k/2]}
 {\sin(\theta_k/2)},
 \qquad \theta_k=\frac{\pi k}{N}.
\end{equation}

Set
\(T_N(m)=\sum_{k=1}^{N-1}\cos(m\theta_k)/\sin(\theta_k/2)\).
Expanding the product of sine squares in
\eqref{eq:supp-b1-exact-diagonal} gives

\[
 \frac14T_N(0)-\frac14T_N(L)-\frac14T_N(a+b)
 +\frac18T_N(2a-1)+\frac18T_N(2b+1).
\]

The zero frequency is
\(T_N(0)=(2N/\pi)\log N+O(N)\).  Discrete cosine orthogonality gives

\begin{equation}\label{eq:supp-b1-parseval}
 \sum_{m=0}^{2N-1}T_N(m)^2
 =N\sum_{k=1}^{N-1}\csc^2\!\left(\frac{\pi k}{2N}\right)
 =\frac{2N(N^2-1)}3.
\end{equation}

Each of the four nonzero frequency maps has at most \(n\) interval
preimages.  Equation \eqref{eq:supp-b1-parseval} therefore makes their
normalized empirical square contribution \(O((\log N)^{-2})\), proving
\eqref{eq:supp-b1-local}.  Since
\(\lvert T_N(m)\rvert\le T_N(0)\) for every integer \(m\), the same exact
diagonal formula also proves \eqref{eq:supp-b1-local-maximum}.
\end{proof}

If \(H_{1,n}^{(D)}=(D_nB_{1,n}B_{1,n}^*D_n)^{1/2}\) and
\(\epsilon_n=1-c_n^n\), the rectangular polar formula gives the pointwise
comparison

\begin{equation}\label{eq:supp-b1-weight-transfer}
 (1-\epsilon_n)^2h_{1,n}(I)
 \le(H_{1,n}^{(D)})_{I,I}
 \le(1-\epsilon_n)^{-1}h_{1,n}(I).
\end{equation}

Indeed, compare \(B_{1,n}^*D_n^2B_{1,n}\) with
\(B_{1,n}^*B_{1,n}\), invert, and apply square-root monotonicity before
taking the polar diagonal.  In a compact critical window,
\(\epsilon_n=O_T((\log n)^{-2})\).  Thus

\begin{equation}\label{eq:supp-physical-b1-local}
 \sup_{0\le t\le T}\frac1{M_n}\sum_I
 \left|\sqrt{M_n}q_n(H_{1,n}^{(D)})_{I,I}
 -\frac{2\sqrt t}{\pi}\right|^2\longrightarrow0.
\end{equation}

For the two-endpoint law, restrict \(Z_n\) to its nonzero proper-interval
rows and write \(H_{2,n}=(Z_nZ_n^*)^{1/2}\).  Map
\((a,b)\) to \((u,v)=(a,n+1-b)\) in \(\Lambda_{n-1}\).

\begin{lemma}[Two-excitation empirical law]
\label{lem:supp-b2-local}
For \(N_n=\binom n2\),

\begin{equation}\label{eq:supp-b2-local}
 \frac1{N_n}\sum_{a<b}
 \left[\frac{(H_{2,n})_{(a,b),(a,b)}}{(\log n)^2}
 -\frac1{\pi^2}\right]^2\longrightarrow0.
\end{equation}
\end{lemma}

\begin{proof}
The interval-poset zeta matrix is invertible; its inverse is the mixed
first-difference operator

\[
 x(a,b)=y(a,b)-y(a+1,b)-y(a,b-1)+y(a+1,b-1).
\]

Let \(D_m=U_m^{-1}\), \(m=n-1\), and let \(V_m\) extend functions on
\(\Lambda_m\) by zero to the full square.  Forward differences preserve the
triangular support, so
\((D_m\otimes D_m)V_m=V_mZ_n^{-1}\).  Since \(x\mapsto x^{-1/2}\) is
operator convex, compression Jensen gives the pointwise comparison

\begin{equation}\label{eq:supp-b2-square-majorant}
 (H_{2,n})_{(a,b),(a,b)}
 \le d_{m,u}d_{m,v},
 \qquad
 d_{m,x}=((U_mU_m^*)^{1/2})_{xx}.
\end{equation}

The sine expansion used in Lemma~\ref{lem:supp-odd-cosine} gives

\[
 \frac1m\sum_x\left(\frac{d_{m,x}}{\log m}-\frac1\pi\right)^2\to0,
 \qquad
 \max_xd_{m,x}\le\frac2\pi\log m+O(1).
\]

The triangular restriction preserves this empirical convergence: for any
sequence \(x_u\) and constant \(a\),

\[
 \frac1{N_n}\sum_{u+v\le m+1}|x_u-a|^2
 \le \frac2m\sum_u|x_u-a|^2.
\]

The product bound and uniform estimate above therefore make the right-hand
side of \eqref{eq:supp-b2-square-majorant} converge to \(1/\pi^2\) in
empirical \(L^2\) over the triangle.  Proposition~\ref{prop:supp-b2-trace} says the
left-hand side has the same empirical mean.  The normalized deficit is
nonnegative and uniformly bounded. If \(\Delta_{a,b}\) denotes this deficit,
then

\[
 \frac1{N_n}\sum_{a<b}\Delta_{a,b}^2
 \le \left(\max_{a<b}\Delta_{a,b}\right)
 \frac1{N_n}\sum_{a<b}\Delta_{a,b}\longrightarrow0.
\]

The triangle inequality in empirical \(L^2\) now proves
\eqref{eq:supp-b2-local}.
\end{proof}

The same polar comparison as in \eqref{eq:supp-b1-weight-transfer}, followed
by the hull order \eqref{eq:supp-hull-order}, transfers
\eqref{eq:supp-b2-local} to all depths \(\ge2\):

\begin{equation}\label{eq:supp-physical-b2-local}
 \sup_{0\le t\le T}\frac1{M_n}\sum_I
 \left|\sqrt{M_n}(\sqrt{A_{\ge2,n}A_{\ge2,n}^*})_{I,I}
 -\frac{\sqrt2\,t}{\pi^2}\right|^2\longrightarrow0.
\end{equation}

\subsection{A Kac--Murdock--Szeg{\H o} diagonal bound}

Let \(T_m(c)=(c^{|j-k|})_{j,k=1}^m\).

\begin{lemma}[Uniform KMS square-root diagonal]
\label{lem:supp-kms-diagonal}
Fix \(C_0<\infty\) and \(\kappa>0\).  There is
\(C=C(C_0,\kappa)\) such that, whenever

\[
 0<c\le1,
 \qquad m(1-c)\le C_0,
 \qquad c^m\ge\kappa,
\]

one has

\begin{equation}\label{eq:supp-kms-bound}
 \max_j(\sqrt{T_m(c)})_{jj}
 \le C\left[m^{-1/2}+\sqrt{1-c}\,\log(m+1)\right].
\end{equation}

For \(c=c_n(t)=1-t/[n(\log n)^2]\), uniformly in
\(0\le t\le T\) and \(1\le m\le n\),

\begin{equation}\label{eq:supp-kms-critical}
 \max_j(\sqrt{T_m(c_n(t))})_{jj}\le\frac{C_T}{\sqrt m}.
\end{equation}
\end{lemma}

\begin{proof}
The cases \(c=1\) and \(m=1\) are immediate.  For \(0<c<1\),

\begin{equation}\label{eq:supp-kms-inverse}
 T_m(c)^{-1}=\frac1{1-c^2}
 \begin{pmatrix}
 1&-c\\
 -c&1+c^2&-c\\
 &\ddots&\ddots&\ddots\\
 &&-c&1
 \end{pmatrix}.
\end{equation}

An eigenvector of the tridiagonal numerator has the form

\begin{equation}\label{eq:supp-kms-eigenvector}
 v_j=A[\sin(j\theta)-c\sin((j-1)\theta)],
 \qquad0<\theta<\pi,
\end{equation}

where \(\theta\) satisfies the right endpoint condition

\begin{equation}\label{eq:supp-kms-endpoint}
 \sin((m+1)\theta)-2c\sin(m\theta)
 +c^2\sin((m-1)\theta)=0.
\end{equation}

Indeed, the left endpoint equation is automatic for
\eqref{eq:supp-kms-eigenvector}, while the right endpoint equation is
equivalent to \(v_{m+1}=cv_m\) after extending the same formula to
\(j=m+1\). Every nonzero
sequence \(y_j=\sin(j\theta+\phi)\), \(1\le j\le m\), satisfies the uniform
sampling inequality

\begin{equation}\label{eq:supp-sampling-inequality}
 \max_j|y_j|^2\le\frac{C_s}{m}\sum_{j=1}^m|y_j|^2.
\end{equation}

To verify it, reduce \(\theta\) to its distance \(d\in[0,\pi/2]\) from
\(\pi\mathbb Z\).  If \(md>2\), the two eigenvalues of the sampling Gram
matrix with rows \((\cos jd,\sin jd)\) are

\[
 \frac m2\mathbin{\pm}\frac{|\sin(md)|}{2\sin d};
\]

the smaller is at least a fixed positive multiple of \(m\).  If \(md\le2\),
write \(u=md\) and use the uniformly nondegenerate basis
\(\cos(ut)\) and \(\sin(ut)/u\), with its continuous limit \(t\) at
\(u=0\).  Its normalized Riemann-sum Gram matrix has a uniformly positive
smallest eigenvalue for \(u\in[0,2]\).  This proves
\eqref{eq:supp-sampling-inequality}, including the finitely many small values
of \(m\) by compactness.  Applied to \eqref{eq:supp-kms-eigenvector}, it gives

\begin{equation}\label{eq:supp-kms-delocalization}
 \max_j|v_j|^2\le C/m
\end{equation}

for every normalized eigenvector.

The standard Cholesky factor of \(T_m(c)\) has first column
\((1,c,\ldots,c^{m-1})^{\mathsf T}\) and lower-right block
\(\sqrt{1-c^2}\,V_{m-1}(c)\), where

\[
 V_s(c)=D_{c,s}U_s^{\mathsf T}D_{c,s}^{-1},
 \qquad D_{c,s}=\operatorname{diag}(1,c,\ldots,c^{s-1}).
\]

The assumption \(c^m\ge\kappa\) bounds the condition number of
\(D_{c,s}\), so \(\|V_s(c)\|_{S_1}\le C\|U_s\|_{S_1}\).  Consequently,

\begin{equation}\label{eq:supp-kms-root-trace}
 \operatorname{tr}\sqrt{T_m(c)}
 \le\sqrt m+C\sqrt{1-c}\,m\log(m+1).
\end{equation}

The spectral resolution, \eqref{eq:supp-kms-delocalization}, and
\eqref{eq:supp-kms-root-trace} prove \eqref{eq:supp-kms-bound}.  In the
critical family, \(c_n(t)^m\ge c_n(t)^n\ge1/2\) for large \(n\), while

\[
 \sqrt{1-c_n(t)}\log(m+1)
 \le\frac{\sqrt T}{\sqrt n}\le\frac{\sqrt T}{\sqrt m\,}.
\]

This proves \eqref{eq:supp-kms-critical} uniformly on \([0,T]\).
\end{proof}

\subsection{Completion of the local law}

Every column of \(B_{1,n}\) is endpoint-additive, because

\[
 \mathbf1_{\{a\le s\le b\}}
 =\mathbf1_{\{a\le s\}}+\mathbf1_{\{s\le b\}}-1.
\]

Lemma~\ref{lem:supp-retained-orthogonality} therefore shows that the three
terms in

\begin{equation}\label{eq:supp-retained-operator}
 K_n=J_{M_n}+q_n^2B_{1,n}Q_{1,n}B_{1,n}^*
 +q_n^4Z_nQ_{2,n}Z_n^*
\end{equation}

have mutually orthogonal row ranges and \(0\preceq K_n\preceq\widehat G_n\).
The projection \(Q_{1,n}\) removes one direction.  Operator monotonicity,
Lemma~\ref{lem:supp-b1-local}, and

\[
 0\le\|B_{1,n}\|_{S_1}-\|B_{1,n}Q_{1,n}\|_{S_1}
 \le\|B_{1,n}\|_{\rm op}=O(n^{3/2})
\]

give the same empirical law for its retained one-excitation diagonal.  In
each two-endpoint cell, the factors in
\eqref{eq:supp-retained-cell-factorization} lose at most two directions.
Lemma~\ref{lem:supp-finite-codimension} applies uniformly because every
block size tends to infinity and
\(\log|E_p|/\log n\to1\).  Products preserve bounded empirical
\(L^2\) convergence.  The omitted same-block cells occupy
\(O(r_n^{-1})\) of the interval triangle.  Hence, with

\[
 \ell_{n,I}=\sqrt{M_n}(\sqrt{K_n})_{I,I},
\]

we obtain, uniformly on every compact parameter interval,

\begin{align}
 \sup_{0\le t\le T}\frac1{M_n}\sum_I
 |\ell_{n,I}-A(t)|^2&\longrightarrow0,
 \label{eq:supp-retained-local}\\
 \sup_{0\le t\le T}\left|
 \frac{\operatorname{tr}\sqrt{K_n}}{\sqrt{M_n}}-A(t)
 \right|&\longrightarrow0.
 \label{eq:supp-retained-trace}
\end{align}

The maximum estimates in Lemmas~\ref{lem:supp-b1-local} and
\ref{lem:supp-finite-codimension}, together with the critical factors
\(q_n\) and \(q_n^2\), also give
\(\sup_{n,t,I}\ell_{n,I}<\infty\) on every compact \(t\)-window.

It remains to show that the nonnegative diagonal excess of
\(\sqrt{\widehat G_n}\) cannot concentrate.  Fix a site \(x\) and consider
all intervals containing it.  Their physical principal Gram matrix is exactly

\begin{equation}\label{eq:supp-anchor-kms-product}
 T_x(c_n)\otimes T_{n+1-x}(c_n).
\end{equation}

Compression Jensen for the operator-concave square root and
Lemma~\ref{lem:supp-kms-diagonal} imply, for every \(I\ni x\),

\begin{equation}\label{eq:supp-anchor-diagonal-bound}
 (\sqrt{G_n})_{I,I}
 \le\frac{C_T}{\sqrt{x(n+1-x)}}.
\end{equation}

Let \(\chi_n(I)=\max_{x\in I}x(n+1-x)\).  Choosing the best anchor gives

\begin{equation}\label{eq:supp-ui-majorant}
 M_n(\sqrt{G_n})_{I,I}^2
 \le C_T\frac{M_n}{\chi_n(I)}.
\end{equation}

The right-hand side is uniformly integrable under the interval prior.  For
large \(K\), only the two edge triangles can satisfy
\(M_n/\chi_n(I)>K\); on the left edge this forces \(b<n/K\), and hence

\begin{equation}\label{eq:supp-ui-tail}
 \frac1{M_n}\sum_I\frac{M_n}{\chi_n(I)}
 \mathbf1_{\{M_n/\chi_n(I)>K\}}
 \le2\sum_{b<n/K}\frac1{n+1-b}
 \le\frac CK.
\end{equation}

Let
\(\widehat\eta_{n,I}=\sqrt{M_n}(\sqrt{\widehat G_n})_{I,I}\). Operator
monotonicity, \eqref{eq:supp-retained-trace}, and
\eqref{eq:supp-unweighted-global-trace-limit} give

\[
 e_{n,I}:=\widehat\eta_{n,I}-\ell_{n,I}\ge0,
 \qquad \frac1{M_n}\sum_Ie_{n,I}\to0.
\]

Put \(\alpha_n=c_n^n\) and
\(\widetilde G_n=D_n\widehat G_nD_n\).  If \(c_n<1\), the interval-zeta
minor makes \(\widehat G_n\) positive definite.  For an invertible square
factor \(\mathsf A_n\mathsf A_n^*=\widehat G_n\),

\[
 \alpha_n^2\mathsf A_n^*\mathsf A_n
 \preceq \mathsf A_n^*D_n^2\mathsf A_n
 \preceq \mathsf A_n^*\mathsf A_n.
\]

Order reversal under inversion, square-root monotonicity, and the polar
formula give the pointwise diagonal inequalities

\begin{equation}\label{eq:supp-physical-diagonal-sandwich}
 \alpha_n^2(\sqrt{\widehat G_n})_{I,I}
 \le(\sqrt{\widetilde G_n})_{I,I}
 \le\alpha_n^{-1}(\sqrt{\widehat G_n})_{I,I}.
\end{equation}

Indeed,

\begin{align*}
 \sqrt{\widehat G_n}
 &=\mathsf A_n(\mathsf A_n^*\mathsf A_n)^{-1/2}\mathsf A_n^*,\\
 \sqrt{\widetilde G_n}
 &=D_n\mathsf A_n
 (\mathsf A_n^*D_n^2\mathsf A_n)^{-1/2}
 \mathsf A_n^*D_n,
\end{align*}

and \(\alpha_n\le(D_n)_{I,I}\le1\).  Since
\(\widetilde G_n\preceq G_n\), equation
\eqref{eq:supp-physical-diagonal-sandwich} yields

\[
 (\sqrt{\widehat G_n})_{I,I}
 \le\alpha_n^{-2}(\sqrt{G_n})_{I,I}.
\]

Thus \eqref{eq:supp-ui-majorant}--\eqref{eq:supp-ui-tail} transfer uniform
integrability to \(\widehat\eta_{n,I}^2\).  If \(c_n=1\), all matrices in
this comparison equal \(J_{M_n}\), so the conclusion is direct.  Splitting
at a fixed level \(R\), using \(e^2\le Re\) below that level and the
uniform-integrability tail above it, yields

\begin{equation}\label{eq:supp-excess-l2}
 \frac1{M_n}\sum_Ie_{n,I}^2\longrightarrow0.
\end{equation}

Combining \eqref{eq:supp-retained-local} and
\eqref{eq:supp-excess-l2} first proves the local law for \(\widehat G_n\).
Because \(\alpha_n=1+O_T((\log n)^{-2})\),
\eqref{eq:supp-physical-diagonal-sandwich} and the bounded empirical second
moments transfer it to \(\widetilde G_n\).  Finally,
Proposition~\ref{prop:supp-hull-compression} gives

\[
 \sqrt{\widetilde G_n}\preceq\sqrt{G_n}
 \preceq\rho_n\sqrt{\widetilde G_n},
 \qquad \rho_n=1+O_T((\log n)^{-2}),
\]

so the same argument transfers the law to the exact physical matrix.

This completes the proof of Theorem~\ref{thm:supp-full-local-law}.

\section{Finite-size numerical diagnostics}
\label{supp:numerics}

\subsection{Safeguarded finite-size SDP endpoints}

Raw solver objectives are retained as diagnostics but are not used directly as
bounds. Let \(\rho_j\) denote the prior weighted canonical states on their
\(d_{\rm supp}\)-dimensional support. For a numerical dual matrix \(Y\), set

\[
\delta_D=\max_j\max\{0,-\lambda_{\min}(Y-\rho_j)\}
\]

and add the recorded roundoff margin \(\zeta_D>0\), scaled to the numerical
problem. Then

\begin{equation}\label{eq:safe-dual-bound}
Y_{\rm safe}=Y+(\delta_D+\zeta_D)I,\qquad
U_{\rm safe}=\operatorname{tr}Y_{\rm safe}
\end{equation}

is rechecked against every dual PSD constraint. For exact matrices satisfying
those constraints, it supplies a feasible upper bound.

For the primal matrices, project each raw \(M_j\) onto the PSD cone and put
\(S=\sum_jM_j^+\). Choose
\(\alpha\geq\max\{1,\lambda_{\max}(S)\}\), with an outward margin chosen for
the numerical scale,
and define

\begin{equation}\label{eq:primal-normalization}
 \begin{aligned}
  Q_j&=\frac{M_j^+}{\alpha},\\
  R&=I-\sum_kQ_k\succeq0,\\
  \overline M_j&=Q_j+\delta_{j1}R.
 \end{aligned}
\end{equation}

This singular-safe construction gives \(\overline M_j\succeq0\) and
\(\sum_j\overline M_j=I\) in exact arithmetic. If \(q\) is the number of
hypotheses, the final positive floor \(f_P>0\), chosen with \(qf_P<1\), is
introduced through the uniform-POVM mixture

\begin{equation}\label{eq:primal-strict-mixture}
 \widehat M_j=(1-qf_P)\overline M_j+f_PI.
\end{equation}

Thus every \(\widehat M_j\) is positive definite and the family remains
complete. The implementation rechecks the resulting POVM. In exact arithmetic,
feasibility and weak duality give

\begin{equation}\label{eq:safe-primal-bound}
L_{\rm safe}=\sum_j\operatorname{tr}(\rho_j\widehat M_j),\qquad
L_{\rm safe}\le P_{\rm opt}\le U_{\rm safe}.
\end{equation}

The implementation uses IEEE double precision and conventional eigensolvers,
with no directed rounding, interval arithmetic, or exact rational
verification. Accordingly, \(L_{\rm safe}\) and \(U_{\rm safe}\) are reported
as safeguarded floating-point endpoints. In exact arithmetic, the sandwich
holds whenever the postprocessed matrices are feasible, and the stored
residuals record the feasibility checks at working precision.

\subsection{Critical-window physical Gram checks}

For each \((n,\tau)\), the diagnostic sets
\(c=1-\tau/[n(\log n)^2]\), construct the full physical Gram matrix explicitly
as a dense matrix, and diagonalize it.
Table~\ref{tab:supp-critical-diagnostics} reports the scaled
trace amplitude, the target \(A(\tau)\), the empirical root-mean-square error
of the scaled diagonal, and the trace loss of the row-orthogonal retained
comparison. The slow finite-size drift, especially at \(\tau=4\), is
consistent with the logarithmic theorem and cautions against interpreting an
initially increasing polar/SRM gap as asymptotic behavior.

\begin{table}[ht]
\centering
\caption{Dense physical Gram diagnostics.  The retained trace gap is
\((\operatorname{tr}\sqrt G-\operatorname{tr}\sqrt K)/\sqrt{M_n}\).}
\label{tab:supp-critical-diagnostics}
\begin{tabular}{rrrrrr}
\toprule
\(\tau\)&\(n\)&\(A(\tau)\)&trace amplitude&local RMS&retained gap\\
\midrule
0.25&12&1.354132&1.442930&0.105208&0.136272\\
0.25&24&1.354132&1.417693&0.084748&0.107393\\
0.25&32&1.354132&1.410576&0.079250&0.099183\\
1.00&12&1.779910&1.965364&0.221155&0.410377\\
1.00&24&1.779910&1.919312&0.186629&0.320009\\
1.00&32&1.779910&1.905585&0.176503&0.295459\\
4.00&12&2.846399&3.165847&0.395940&1.490082\\
4.00&24&2.846399&3.124405&0.388835&1.105722\\
4.00&32&2.846399&3.107347&0.381736&1.012978\\
\bottomrule
\end{tabular}
\end{table}

At \(n=32\), the differences between the full square-root diagonal and the
sum of the sectorwise square-root diagonals have the ranges

\[
\begin{array}{c|ccc}
\tau&0.25&1&4\\ \hline
\text{range}&[-0.2175,0.1660]&[-0.5349,0.4288]&[-1.0370,1.1270].
\end{array}
\]

Both signs occur at the working precision, providing numerical evidence
against coordinatewise square-root additivity. The minimum eigenvalue of
\(\widehat G_n-K_n\) lay
between \(-2.9\times10^{-15}\) and \(-2.2\times10^{-15}\), and the retained
row-orthogonality residuals were below \(1.8\times10^{-13}\), consistent with
positive semidefiniteness at the working precision.
